% exposition: G
% scope: линия G — уравнение -y''(x)-a(x)y''(-x)=lambda y с общими двухточечными условиями, без младших членов; записи C24, C25, C26. Не входят: линия D (базисность при условиях Дирихле), идеи, базисность, случаи p=0 и p<0
\documentclass[12pt,a4paper]{article}
\usepackage[utf8]{inputenc}
\usepackage[T2A]{fontenc}
\usepackage[russian]{babel}
\usepackage{amsmath,amssymb,amsthm}
\usepackage{longtable,array}
\usepackage[margin=2.5cm]{geometry}
\usepackage[unicode,hidelinks]{hyperref}
\usepackage{indentfirst}
\setlength{\emergencystretch}{3em}
\numberwithin{equation}{section}
\newtheorem{theorem}{Теорема}
\newtheorem{lemma}{Лемма}
\newtheorem{proposition}{Предложение}
\newtheorem{corollary}{Следствие}
\theoremstyle{definition}
\newtheorem{definition}{Определение}
\newtheorem{remark}{Замечание}
\renewcommand{\qedsymbol}{}
\newcommand{\step}[1]{\smallskip\noindent\textbf{Шаг #1.}}
\newcommand{\R}{\mathbb R}
\let\C\relax\newcommand{\C}{\mathbb C}
\newcommand{\norm}[1]{\lVert#1\rVert}
\newcommand{\abs}[1]{\lvert#1\rvert}
\newcommand{\Imm}{\operatorname{Im}}
\newcommand{\Ree}{\operatorname{Re}}
\newcommand{\Uf}{\mathfrak U}
\newcommand{\Mm}{\mathcal N}
\newcommand{\Dd}{\mathbf D}
\newcommand{\Zf}{\mathcal Z_F}
\title{Регулярные краевые задачи для уравнения второго порядка\\
с отражением аргумента в старшей производной}
\author{Л.\,И.~Дядченко}
\date{}
\begin{document}
\maketitle
\begin{abstract}
Рассматривается спектральная краевая задача с отражением аргумента в старшей производной
$-y''(x)-a(x)y''(-x)=\lambda y(x)$ на $[-1,1]$ с общими двухточечными условиями.
Вещественный коэффициент $a\in W^1_1(-1,1)$ отделён от нуля и удовлетворяет неравенству
$a(x)a(-x)\le p_0<1$. Даны определения регулярной и сильно регулярной задач.
Для регулярной задачи получена асимптотика $\sqrt{\lambda_k}$ с остатком $o(1)$;
для сильно регулярной "--- также равномерная асимптотика собственных функций и
производных. При $a\in W^2_1(-1,1)$ и сильной регулярности
остатки имеют порядок $O(1/k)$. Результаты продемонстрированы
на задаче Дирихле и на задаче с условиями $y(1)=y(-1)$, $y'(1)=-y'(-1)$.
\end{abstract}

\section{Введение}

Рассматривается спектральная краевая задача с отражением аргумента в старшей производной
\begin{gather}
 -y''(x)-a(x)\,y''(-x)=\lambda y(x),\qquad x\in[-1,1],\label{eq:eq}\\
 U_j(y)=\alpha_{j1}y'(-1)+\alpha_{j0}y(-1)+\beta_{j1}y'(1)+\beta_{j0}y(1)=0,\qquad j=1,2,\label{eq:bc}
\end{gather}
с комплексными коэффициентами $\alpha_{ji}$, $\beta_{ji}$ и линейно независимыми формами
$U_1,U_2$. Функция $a$ вещественна и удовлетворяет условиям
\[
 a\in W^1_1(-1,1),\qquad \abs{a(x)}\ge m>0,\qquad a(x)a(-x)\le p_0<1.
\]

При $a\equiv0$ уравнение \eqref{eq:eq} превращается в $-y''=\lambda y$, то есть в
классическую задачу; это предельный случай, в котором нарушено условие $\abs a\ge m$.

Для уравнений второго порядка с отражением в старшей части известны результаты о
свойствах систем корневых функций. А.\,А.~Копжасарова и А.~Сарсенби~\cite{KopSar}
для уравнения с постоянными комплексными коэффициентами и общими двухточечными условиями
получили условия базисности Рисса. А.~Сарсенби и А.\,М.~Сарсенби~\cite{Sars} ввели функцию
Грина для уравнения с постоянным коэффициентом при $y''(-x)$ и суммируемым потенциалом,
установили равносходимость спектральных разложений при условиях $y'(-1)+b_1y(-1)=0$,
$y'(1)+b_2y(1)=0$ и получили достаточные условия базисности собственных функций.
Л.\,В.~Крицков и А.\,М.~Сарсенби~\cite{KriSar} для оператора с постоянным коэффициентом
при $y''(-x)$, суммируемыми потенциалами и произвольной абсолютно непрерывной инволюцией
нашли необходимые и достаточные условия безусловной базисности. Л.\,В.~Крицков,
М.\,А.~Садыбеков и А.\,М.~Сарсенби~\cite{KSadSar} для нелокальных условий
$u(-1)=\beta u(1)$, $u'(-1)=u'(1)$ и коэффициента $\alpha$ показали, что свойства системы определяются числом
$r=\sqrt{(1-\alpha)/(1+\alpha)}$: при иррациональном $r$ система полна и минимальна,
но не образует базиса, при рациональном "--- образует безусловный базис.

Обратная спектральная задача для уравнения с отражением решена
Н.\,П.~Бондаренко~\cite{Bond}: сведение к матричной задаче Штурма--Лиувилля при условиях
Дирихле и постоянном коэффициенте при $y''(-x)$ позволило восстановить потенциалы по
спектру. А.~Бейсебаева и А.\,М.~Сарсенби~\cite{BeiSar} для уравнения
$-u''(-x)+q(x)u=\lambda u$ с чётным комплекснозначным $q$ и условиями $u(-1)=0$,
$u'(1)=u'(-1)$ нашли асимптотики собственных значений, собственных функций и их
производных с остатком $O(1/k)$. Прямая вторая производная в это уравнение не входит,
поэтому главные части обеих серий выписываются через синус и косинус, а асимптотики
получаются сведением к интегральным уравнениям Вольтерра. Порядок остатка $O(1/k)$ тот же,
что в теореме~\ref{thm:o1k}; отличие в том, что в~\cite{BeiSar} главная часть постоянна,
а краевые условия образуют одну пару форм. Во всех перечисленных работах коэффициенты при
вторых производных постоянны.

Постановка восходит к теории Биркгофа. Хорошо известно (например,~\cite{Na}),
что регулярность краевых условий по Биркгофу обеспечивает асимптотическое описание
собственных значений по главной части характеристического определителя, а сильная
регулярность "--- асимптотику отдельных собственных функций. Для операторов с инволюцией $L=JP+Q$,
где $P$ и $Q$ "--- обыкновенные дифференциальные выражения порядков $n$ и $m$, а $J$
"--- отражение, определения регулярной задачи даны В.\,Е.~Владыкиной и
А.\,А.~Шкаликовым~\cite{VS}. Основное внимание там уделено случаю $n>m$; случай же $n=m$
не рассмотрен~\cite[с.~644]{VS}. В уравнении \eqref{eq:eq} обе вторые производные входят в
старшую часть, то есть в обозначениях~\cite{VS} это случай $n=m=2$, и определения
из~\cite{VS} непосредственно к задаче \eqref{eq:eq}, \eqref{eq:bc} не относятся.

Различие между постоянным и переменным коэффициентом видно при переходе к системе
на $[0,1]$ с заменой $w(x)=(y(x),y(-x))^\top$:
\[
 -M(x)w''=\lambda w,\qquad M(x)=\begin{pmatrix}1&a(x)\\a(-x)&1\end{pmatrix}.
\]
При постоянном $a$ эта система распадается на уравнения для чётной и нечётной частей
решения. При переменном $a$ собственные векторы её матрицы, вообще говоря, зависят от $x$;
переход к ним порождает дополнительные члены, и вместо двух независимых скалярных
уравнений получается система с переменными коэффициентами. Условия на $a$ обеспечивают
положительность и разделённость собственных значений её главной матрицы $M(x)$,
\[
 \sigma_\pm(x)=1\pm\sqrt{a(x)a(-x)},
\]
которым отвечают фазовые интегралы и фазовые длины
\[
 t_\pm(x)=\int_0^x\frac{ds}{\sqrt{\sigma_\pm(s)}},\qquad \ell_\pm=t_\pm(1),
\]
входящие в главные части решений.

Для систем первого порядка с переменной главной матрицей определения регулярной и сильно
регулярной задач предложены А.\,А.~Шкаликовым~\cite[с.~798]{Sh}:
регулярность определяется отличием от нуля коэффициентов при крайних экспонентах
характеристического определителя упрощённой задачи, сильная регулярность требует
дополнительно асимптотической простоты и равномерной отделённости его нулей. При сильной
регулярности собственные функции асимптотически совпадают с собственными функциями
упрощённой задачи~\cite[теорема~2]{Sh}. Для операторов первого порядка
$\alpha(x)y'(x)+y'(-x)+q(x)y(x)+r(x)y(-x)$ с переменным $\alpha(x)$ и одним краевым
условием $y(-1)=hy(1)$ получены условия базисности корневых функций~\cite{GrSh},
\cite{ShInv}. Методы, развитые там, не удалось распространить на уравнение второго порядка.

Перенос такого подхода к регулярности на задачу \eqref{eq:eq}, \eqref{eq:bc} даёт определитель главных
членов краевых форм. Он вычисляется явно через коэффициенты $U_j$, значения $a(\pm1)$ и
фазовые длины $\ell_\pm$, а после выделения степени параметра $\rho=\sqrt\lambda$
представляет собой сумму четырёх экспонент с частотами $\pm(\ell_++\ell_-)$ и
$\pm(\ell_+-\ell_-)$; эту сумму обозначаем через $F(\rho)$. Задача называется регулярной,
если коэффициенты при крайних экспонентах в $F$ отличны от нуля, и сильно регулярной,
если, кроме того, большие нули $F$ просты и равномерно отделены друг от друга
(определение~\ref{def:reg}).

Цель настоящей работы "--- дать эти определения для уравнения с отражением в старшей
производной и доказать для обоих классов асимптотические формулы. Основной результат
получен в теореме~\ref{thm:eig}: для регулярной задачи при подходящей нумерации
\[
 \sqrt{\lambda_k}=\rho_k^0+o(1),\qquad k\to\infty,
\]
где $\rho_k^0$ "--- нули $F$ с положительной вещественной частью при больших $k$,
а нумерация учитывает кратности. Для сильно регулярной задачи получены равномерные
асимптотические формулы для собственных функций и их производных с остатком $o(1)$
после масштабирования производной (теорема~\ref{thm:ef}), а при дополнительном условии
$a\in W^2_1(-1,1)$ остатки в асимптотиках $\sqrt{\lambda_k}$ и собственных функций имеют
порядок $O(1/k)$ (теорема~\ref{thm:o1k}). На примере задачи Дирихле продемонстрированы введённые определения
и получен критерий сильной регулярности через отношение
фазовых длин (предложение~\ref{prop:dir}); достаточные условия сильной регулярности для
краевых условий вида $(1,0)$ (см. п.~\ref{ss:forms}) даны в предложении~\ref{prop:10}, а пример нерегулярной задачи, у которой при
чётном $a$ каждое комплексное число является собственным значением, "--- в
предложении~\ref{prop:refl}.

Асимптотические представления решений получены с помощью теорем А.\,П.~Косарева и
А.\,А.~Шкаликова для систем с большим параметром~\cite[теорема~1.5]{KS1},
\cite[теорема~1.2]{KS2}. Переход от этих представлений к спектральным асимптотикам
проводится для краевых форм $U_j$, включающих члены нулевого порядка в условиях.
При переходе к системе первого порядка такие члены дают зависимость краевых матриц от
параметра, отсутствующую в постановке~\cite{Sh}, но учитываются как малая поправка и
потому не меняют ни определения регулярности, ни главных частей асимптотических формул.

\section{Обозначения и основные результаты}

\subsection{Условия на коэффициент}\label{ss:H}

Всюду далее функция $a$ вещественна и при некоторых $m>0$, $p_0<1$ удовлетворяет условиям
\begin{equation}\label{eq:H}
 \text{(i)}\ a\in W^1_1(-1,1);\qquad
 \text{(ii)}\ \abs{a(x)}\ge m;\qquad
 \text{(iii)}\ p(x):=a(x)a(-x)\le p_0,\quad x\in[-1,1].
\end{equation}
По (i) функция $a$ непрерывна, а по (ii) не обращается в нуль и потому знакопостоянна;
положим $\varepsilon=\operatorname{sign}a$. Тогда $m^2\le p(x)\le p_0$, и собственные значения
матрицы $M(x)$ из введения
\[
 \sigma_\pm(x)=1\pm\sqrt{p(x)}
\]
положительны и различны равномерно по $x$:
\[
 \sigma_-\ge1-\sqrt{p_0}>0,\qquad \sigma_+-\sigma_-=2\sqrt p\ge2m.
\]
Первое неравенство позволяет обратить главную матрицу системы, второе "--- разделить
два её собственных значения при построении асимптотических решений.

Число $\lambda\in\C$ называем \emph{собственным значением} задачи \eqref{eq:eq}, \eqref{eq:bc},
если существует функция $y\in W^2_1(-1,1)$, $y\not\equiv0$, которая удовлетворяет \eqref{eq:eq}
почти всюду и условиям \eqref{eq:bc}; такая $y$ "--- \emph{собственная функция}. В \S\,3
построена целая функция $\Delta$ (характеристический определитель), нули которой "--- в точности
собственные значения (лемма~\ref{lem:fold}); \emph{кратностью} собственного значения называем
кратность нуля $\Delta$ (она может не совпадать с алгебраической кратностью собственного
значения). 
%Через $C$ и $c$ без индексов обозначаем положительные постоянные, не зависящие от
%$\rho$, $k$ и $x$; в разных местах они разные.

С задачей \eqref{eq:eq}, \eqref{eq:bc} свяжем оператор $L_a$ в $L_2(-1,1)$ комплекснозначных функций:
\[
 (L_ay)(x)=-y''(x)-a(x)y''(-x),\qquad
 D(L_a)=\{y\in W_2^2(-1,1):\ U_1(y)=U_2(y)=0\}.
\]
При условии $a(x)a(-x)\le p_0<1$ всякое решение уравнения \eqref{eq:eq} из $W^2_1(-1,1)$
лежит в $W_2^2(-1,1)$: исключая $y''(-x)$ из уравнения и его отражения, получаем
\[
 \bigl(1-a(x)a(-x)\bigr)y''(x)=\lambda\bigl(a(x)y(-x)-y(x)\bigr),
\]
а множитель $1-a(x)a(-x)$ отделён от нуля. Поэтому собственные функции задачи и оператора
$L_a$ совпадают.

Оператор $L_a$, вообще говоря, несамосопряжён. При переменном $a$ уже дифференциальное
выражение не является формально самосопряжённым: если $Rf=f(-\cdot)$, то $(aR)^*=Ra$,
причём $(Raf)(x)=a(-x)f(-x)$, и равенство $aR=Ra$ выполнено только для чётного $a$.
При чётном $a$, в частности при постоянном, выражение формально самосопряжено, но краевые
условия \eqref{eq:bc} с комплексными коэффициентами могут не быть самосопряжёнными. Пример
предложения~\ref{prop:refl} показывает, что и при формально самосопряжённом выражении
вырожденные краевые условия меняют картину: условия \eqref{eq:refl} автоматически выполнены
на чётных функциях, и каждое $\lambda\in\C$ становится собственным значением.

\subsection{Фазовые интегралы}\label{ss:chan}

Пусть $\lambda=\rho^2$ и $w(x)=(y(x),y(-x))$, $x\in[0,1]$; тогда $-M(x)w''=\lambda w$
(лемма~\ref{lem:fold}). Матрица $M(x)$ имеет собственные векторы
\[
 \mathbf u_\pm(x)=(1,\pm\nu(x))^\top,\qquad \nu(x)=\varepsilon\sqrt{a(-x)/a(x)},
\]
отвечающие собственным значениям $\sigma_\pm(x)=1\pm\sqrt{a(x)a(-x)}$. При постоянной матрице 
$M$ переход к собственному базису даёт два скалярных уравнения с решениями $e^{\pm i\rho x/\sqrt{\sigma_\pm}}$.
Для переменной матрицы положим $\tau_\pm=\sigma_\pm^{-1/2}$; соответствующие
экспоненциальные множители в главных частях решений имеют вид $e^{\pm i\rho t_\pm(x)}$, где
\begin{equation}\label{eq:not}
 t_\pm(x)=\int_0^x\tau_\pm,\qquad \ell_\pm=t_\pm(1),\qquad L=\ell_++\ell_-.
\end{equation}
Функции $t_\pm$ назовём \emph{фазовыми интегралами}, а числа $\ell_\pm$ "--- \emph{фазовыми
длинами}. Асимптотическое представление фундаментальной матрицы с введенными множителями дано в
лемме~\ref{lem:fsr}. Поскольку $\sigma_+>\sigma_-$, имеем $\tau_+<\tau_-$, откуда $\ell_+<\ell_-$.
Функции $\tau_\pm$ и $\nu$ абсолютно непрерывны, так как в их формулах знаменатели и
подкоренные выражения отделены от нуля. Поэтому производные матрицы собственных
векторов суммируемы, что и позволяет применить теоремы об асимптотике систем.

\subsection{Главные члены краевых форм}\label{ss:forms}

Формы $U_1,U_2$ предполагаются линейно независимыми, в частности ненулевыми. \emph{Порядок}
$\kappa_j\in\{0,1\}$ формы $U_j$ "--- наибольший порядок производной, входящей в неё с ненулевым
коэффициентом. Краевые условия не меняются при замене пары форм парой
$(c_{11}U_1+c_{12}U_2,\,c_{21}U_1+c_{22}U_2)$ с невырожденной матрицей $(c_{ij})$, поэтому пару
можно выбрать так, чтобы сумма порядков $\kappa_1+\kappa_2$ была наименьшей. Такую пару назовём
\emph{нормированной}; нормировка "--- это выбор пары, а не дополнительное свойство задачи.
Всюду далее через $U_1,U_2$ обозначена нормированная пара. С точностью до перестановки форм возможны три
\emph{вида} $(\kappa_1,\kappa_2)$: $(0,0)$ "--- производные не входят ни в одно условие, $(1,1)$ "---
входят в оба, и $(1,0)$ "--- входят только в $U_1$.

Для вида $(0,0)$ выбор форм ничего не меняет. Для вида $(1,1)$ нормированность нарушается ровно
тогда, когда векторы коэффициентов при производных $(\alpha_{11},\beta_{11})$ и
$(\alpha_{21},\beta_{21})$ пропорциональны: тогда подходящая комбинация форм не содержит
производных. Условие нормированности равносильно неравенству
$\beta_{11}\alpha_{21}-\alpha_{11}\beta_{21}\ne0$: при его выполнении производные входят в обе
формы любой пары, полученной невырожденной заменой. Вид $(1,0)$ нормирован всегда.

Пара $(y(x),y(-x))$ должна удовлетворять в нуле двум условиям:
$w_1(0)=w_2(0)$ и $w_1'(0)+w_2'(0)=0$. Они выражают совпадение значений и
односторонних производных исходной функции; далее называем их \emph{условиями склейки}.
В нуле собственные векторы равны $\mathbf u_\pm(0)=(1,\pm\varepsilon)^\top$.
Для вектора $(1,1)^\top$ условие на значения выполнено тождественно, а условие на
производные выбирает косинус из пары экспонент. Для вектора $(1,-1)^\top$ условие на
производные выполнено тождественно, а условие на значения выбирает синус. Поэтому
положим $\chi_+=\cos$, $\chi_-=\sin$ при $\varepsilon=1$ и
$\chi_+=\sin$, $\chi_-=\cos$ при $\varepsilon=-1$.

Остаются две главные части решений, с тригонометрическими множителями
$\chi_+(\rho t_+)$ и $\chi_-(\rho t_-)$. Амплитуды этих частей нормируем значением $1$
в точке $x=1$; их вид на всём отрезке будет нужен для собственных функций
(формула~\eqref{eq:Psi}). Для вычисления краевых форм достаточно значений в концах.
В главной части, отвечающей $\sigma_\pm$, значения $(y(1),y(-1))$ пропорциональны
$(1,\pm\nu(1))$, а производные $(y'(1),y'(-1))$ "--- $(1,\mp\nu(1))$.
Главный член производной содержит множитель
$\rho\tau_\pm(1)\chi_\pm'(\rho\ell_\pm)$. Здесь и далее штрих у $\chi_\pm$ означает
производную по аргументу.

Таким образом, для формы порядка~$1$ члены со значениями функции на одну степень $\rho$
ниже членов с производными. Разделив строку формы $U_j$ на $\rho^{\kappa_j}$ и оставив
главные члены, получаем матрицу $\Mm(\rho)$ с двумя столбцами. Её строка $j$ равна
\begin{equation}\label{eq:M}
\begin{aligned}
&\bigl((\beta_{j0}+\nu(1)\alpha_{j0})\chi_+(\rho\ell_+),\
 (\beta_{j0}-\nu(1)\alpha_{j0})\chi_-(\rho\ell_-)\bigr) &&\text{при } \kappa_j=0,\\
&\bigl(\tau_+(1)(\beta_{j1}-\nu(1)\alpha_{j1})\chi_+'(\rho\ell_+),\
 \tau_-(1)(\beta_{j1}+\nu(1)\alpha_{j1})\chi_-'(\rho\ell_-)\bigr) &&\text{при } \kappa_j=1.
\end{aligned}
\end{equation}
Её ядро задаёт коэффициенты линейной комбинации главных частей, удовлетворяющей
главным членам обоих краевых условий. Поэтому нули её определителя определяют
приближённые значения параметра $\rho$. Выделяя ненулевой постоянный множитель из
$\det\Mm(\rho)$, получаем функцию
\begin{equation}\label{eq:F}
 F(\rho)=\begin{cases}
 \chi_+(\rho\ell_+)\,\chi_-(\rho\ell_-) & \text{для вида } (0,0),\\
 -\chi_+'(\rho\ell_+)\,\chi_-'(\rho\ell_-) & \text{для вида } (1,1),\\
 \varepsilon\bigl(\mathcal B\chi_+(\rho\ell_+)\,\chi_-'(\rho\ell_-)
 -\mathcal A\chi_+'(\rho\ell_+)\,\chi_-(\rho\ell_-)\bigr) & \text{для вида } (1,0),
 \end{cases}
\end{equation}
где для вида $(1,0)$
\[
 \mathcal A=\tau_+(1)\bigl(\beta_{11}-\nu(1)\alpha_{11}\bigr)\bigl(\beta_{20}-\nu(1)\alpha_{20}\bigr),
 \qquad
 \mathcal B=\tau_-(1)\bigl(\beta_{11}+\nu(1)\alpha_{11}\bigr)\bigl(\beta_{20}+\nu(1)\alpha_{20}\bigr).
\]
Равенство $\det\Mm=cF$ с постоянной $c\ne0$ проверяется в
лемме~\ref{lem:kernel}. Для определения регулярности выделим крайние экспоненты
в $F$. Положим $d=\ell_+-\ell_-$; тогда $0<\abs d<L$. По формулам произведений
синусов и косинусов при обоих значениях $\varepsilon$ получаем
\begin{equation}\label{eq:Ftrig}
 F(\rho)=\begin{cases}
 \frac12\bigl(\sin\rho L-\varepsilon\sin\rho d\bigr) & \text{для вида } (0,0),\\
 \frac12\bigl(\sin\rho L+\varepsilon\sin\rho d\bigr) & \text{для вида } (1,1),\\
 \frac12\bigl(\varepsilon(\mathcal B-\mathcal A)\cos\rho L+(\mathcal A+\mathcal B)\cos\rho d\bigr)
 & \text{для вида } (1,0).
 \end{cases}
\end{equation}
Следовательно, $F$ имеет представление
\begin{equation}\label{eq:Fexp}
 F(\rho)=\theta e^{i\rho L}+\theta'e^{-i\rho L}+\theta_de^{i\rho d}+\theta_d'e^{-i\rho d},
 \qquad \abs{\theta'}=\abs\theta.
\end{equation}
Частоты различны, поэтому коэффициенты определены однозначно. Равенство модулей
крайних коэффициентов следует из \eqref{eq:Ftrig}: слагаемое с частотой $L$
пропорционально либо синусу, либо косинусу. В частности,
\[
 \theta=-\frac i4\ \text{для видов } (0,0),\ (1,1),\qquad
 \theta=\frac{\varepsilon(\mathcal B-\mathcal A)}4\ \text{для вида } (1,0).
\]
\subsection{Регулярность}\label{ss:model}

Связь $F$ с исходной задачей даёт асимптотика характеристического определителя
(лемма~\ref{lem:det}):
\[
 \rho^{1-\kappa_1-\kappa_2}\Delta(\rho^2)=c_\ast\bigl(F(\rho)+\eta(\rho)\bigr),\qquad
 \sup_{\abs\rho\ge r}\abs{\eta(\rho)}e^{-L\abs{\Imm\rho}}\to0\quad(r\to\infty),
\]
где $c_\ast\ne0$. Остаток $\eta$ мал по сравнению с весом $e^{L\abs{\Imm\rho}}$,
который служит естественным масштабом $F$: каждое слагаемое в \eqref{eq:Fexp}
есть экспонента с частотой не больше $L$. Поэтому при $\theta\ne0$ и $\abs d<L$
одна из крайних экспонент преобладает над остальными вне горизонтальной полосы.
Там $F$ оценивается снизу величиной $ce^{L\abs{\Imm\rho}}$, и малая поправка
$\eta$ не создаёт больших нулей $\Delta(\rho^2)$. В полосе сравнение нулей
проводится по теореме Руше на контурах, отстоящих от нулей $F$
(доказательство теоремы~\ref{thm:eig}, лемма~\ref{lem:F}). Для асимптотики
отдельных собственных функций требуется также устойчивость ядра матрицы
\eqref{eq:M}; в доказательстве её обеспечивает простота и равномерная
отделённость больших нулей $F$.

\begin{definition}\label{def:reg}
Задачу \eqref{eq:eq}, \eqref{eq:bc} назовём \emph{регулярной}, если коэффициент при крайней
экспоненте $e^{i\rho L}$ в \eqref{eq:Fexp} отличен от нуля, то есть $\theta\ne0$. Регулярную задачу
назовём \emph{сильно регулярной}, если все нули $F$ из полуплоскости $\Ree\rho\ge0$, кроме конечного числа, простые,
и существует $\delta_s>0$, при котором $\abs{\rho^0-\rho^{0\prime}}\ge\delta_s$ для любых двух
различных нулей $\rho^0,\rho^{0\prime}$ функции $F$ из этой полуплоскости.
\end{definition}

Для видов $(0,0)$ и $(1,1)$ регулярность выполнена всегда: $\theta=-i/4$. Для вида $(1,0)$ условие
$\theta\ne0$ равносильно неравенству $\mathcal A\ne\mathcal B$, в которое входят $\tau_\pm(1)$ и
$\nu(1)$, то есть значения $a(\pm1)$; поэтому мы говорим о регулярности задачи, а не условий.

По \eqref{eq:Ftrig} функция $F$ нечётна для видов $(0,0)$, $(1,1)$ и чётна для
вида $(1,0)$; множество её нулей обозначим через $\Zf$, тогда $\Zf=-\Zf$.
В определении достаточно полуплоскости $\Ree\rho\ge0$, поскольку точки $\rho$ и
$-\rho$ отвечают одному значению $\lambda=\rho^2$. При регулярности нули $F$ лежат
в горизонтальной полосе (лемма~\ref{lem:F}(a)), так что большие нули правой
полуплоскости далеки от левой. Их простота и отделённость влекут те же свойства
для больших нулей во всей плоскости. Через $r_s$ обозначаем число, при котором
нули сильно регулярной задачи с $\Ree\rho\ge0$ и $\abs\rho\ge r_s$ просты.

Для задачи Дирихле при $a>0$ получаем $F(\rho)=\cos(\rho\ell_+)\sin(\rho\ell_-)$.
Крайние коэффициенты ненулевые, поэтому задача регулярна. Однако две серии
нулей могут пересекаться или сближаться: сильная регулярность требует их
равномерной отделённости. Точный критерий дан в предложении~\ref{prop:dir}.
Пример нерегулярной задачи, для которой $F\equiv0$, приведён в
предложении~\ref{prop:refl}.

\subsection{Собственные значения}\label{ss:eig}

Для сопоставления нулей в переменной $\lambda$, включая нуль, восстановим степень
$\rho$, выделенную из характеристического определителя. Положим
\begin{equation}\label{eq:G}
 G(\rho^2)=\rho^{\kappa_1+\kappa_2-1}F(\rho).
\end{equation}
Правая часть "--- чётная целая функция: для видов $(0,0)$ и $(1,1)$ функция $F$
нечётна, показатель равен соответственно $-1$ и $1$, причём $F(0)=0$;
для вида $(1,0)$ функция $F$ чётна и показатель равен нулю. Поэтому равенство
\eqref{eq:G} задаёт единственную целую функцию $G$ переменной $\lambda$.
При $\rho\ne0$ кратность нуля $G$ в точке $\lambda=\rho^2$ равна кратности нуля
$F$ в точке $\rho$; в нуле множитель в \eqref{eq:G} учитывается отдельно.

\begin{theorem}\label{thm:eig}
Пусть $a$ удовлетворяет условиям \eqref{eq:H}, формы $U_1$, $U_2$ линейно независимы,
а пара $U_1,U_2$ нормирована, и задача \eqref{eq:eq}, \eqref{eq:bc} регулярна. Тогда $\Delta\not\equiv0$ и:
\begin{enumerate}
\item[(a)] существует $h>0$, при котором каждое собственное значение $\lambda$ с достаточно
 большим $\abs\lambda$ удовлетворяет неравенству $\abs{\Imm\sqrt\lambda}\le h$;
\item[(b)] нули функций $\Delta$ и $G$, каждый повторённый столько раз, какова его кратность,
 можно занумеровать как $\lambda_k$, $\lambda_k^0$, $k\ge1$, так, что
 $\rho_k-\rho_k^0\to0$ при $k\to\infty$; здесь $\rho_k$ и $\rho_k^0$ "--- квадратные корни из
 $\lambda_k$ и $\lambda_k^0$ с неотрицательной вещественной частью. При больших $k$ вещественные
 части $\rho_k$, $\rho_k^0$ положительны, так что корни определены однозначно, а $\rho_k^0$ "---
 нуль $F$, кратность которого равна кратности нуля $\lambda_k^0$ функции $G$;
\item[(c)] если задача сильно регулярна, то при больших $k$ нули $\lambda_k$ функции $\Delta$
 простые, а собственные подпространства одномерны.
\end{enumerate}
\end{theorem}

\subsection{Главная часть собственной функции}\label{ss:ef}

Нуль $F$ задаёт приближённое собственное значение, но не коэффициенты соответствующей
собственной функции. Эти коэффициенты определяются ядром матрицы $\Mm$ из
\eqref{eq:M}. Чтобы выписать саму главную часть на всём отрезке, к
тригонометрическим множителям нужно добавить амплитуды. Положим
\[
 D_\pm(x)=\Bigl(\frac{\abs{\nu(1)}\,\tau_\pm(1)}{\abs{\nu(x)}\,\tau_\pm(x)}\Bigr)^{1/2}.
\]
Они нормированы условиями $D_\pm(1)=1$ и удовлетворяют равенствам
$D_\pm'/D_\pm=-\nu'/(2\nu)-\tau_\pm'/(2\tau_\pm)$ почти всюду. Эти равенства
получаются при выделении диагональной части системы первого порядка в \S\,3.
Главные части пары значений и пары масштабированных производных запишем через
два коэффициента $\xi$, $\zeta$:
\begin{equation}\label{eq:Psi}
\begin{gathered}
 \begin{pmatrix}y(x)\\ y(-x)\end{pmatrix}
 =\xi D_+(x)\chi_+(\rho t_+(x))\,\mathbf u_+(x)
 +\zeta D_-(x)\chi_-(\rho t_-(x))\,\mathbf u_-(x),\\
 \frac1\rho\begin{pmatrix}y'(x)\\ -y'(-x)\end{pmatrix}
 =\xi D_+(x)\tau_+(x)\chi_+'(\rho t_+(x))\,\mathbf u_+(x)
 +\zeta D_-(x)\tau_-(x)\chi_-'(\rho t_-(x))\,\mathbf u_-(x),
\end{gathered}
\end{equation}
Здесь $x\in[0,1]$ и $\rho\ne0$. Две строки \eqref{eq:Psi} задают главные
члены отдельно: вторая не получается точным дифференцированием первой, поскольку
производные амплитуд и собственных векторов входят в остаток. Условия склейки
для этих главных членов выполнены при любых $\xi$, $\zeta$
(лемма~\ref{lem:kernel}). После подстановки в главные члены краевых форм получаем
$\Mm(\rho)(\xi,\zeta)^\top=0$. В простом нуле $\rho_k^0$ функции $F$ это ядро
одномерно, а сильная регулярность позволяет перенести его на точную матрицу
краевых условий с ошибкой $o(1)$.

\begin{theorem}\label{thm:ef}
Пусть выполнены условия теоремы~\ref{thm:eig}, задача \eqref{eq:eq}, \eqref{eq:bc} сильно
регулярна, и $\lambda_k$, $\lambda_k^0$, $\rho_k$, $\rho_k^0$ "--- нумерация теоремы~\ref{thm:eig}.
Тогда существует $k_0$, при котором для $k\ge k_0$:
\begin{enumerate}
\item[(a)] $\lambda_k$ "--- простой нуль $\Delta$, и собственное подпространство одномерно;
\item[(b)] ранг матрицы $\Mm(\rho_k^0)$ равен единице; пусть $(\xi_k,\zeta_k)$,
 $\abs{\xi_k}^2+\abs{\zeta_k}^2=1$, порождает её ядро;
\item[(c)] существует собственная функция $y_k$, отвечающая $\lambda_k$, для которой равномерно по
 $x\in[0,1]$
 \begin{equation}\label{eq:efas}
 \begin{gathered}
  \begin{pmatrix}y_k(x)\\ y_k(-x)\end{pmatrix}
  =\xi_kD_+(x)\chi_+(\rho_k^0t_+(x))\,\mathbf u_+(x)\\
  +\zeta_kD_-(x)\chi_-(\rho_k^0t_-(x))\,\mathbf u_-(x)+o(1),\\
  \frac1{\rho_k}\begin{pmatrix}y_k'(x)\\ -y_k'(-x)\end{pmatrix}
  =\xi_kD_+(x)\tau_+(x)\chi_+'(\rho_k^0t_+(x))\,\mathbf u_+(x)\\
  +\zeta_kD_-(x)\tau_-(x)\chi_-'(\rho_k^0t_-(x))\,\mathbf u_-(x)+o(1),
 \end{gathered}
 \end{equation}
\end{enumerate}
При этом главная часть в \eqref{eq:efas} не мала: при $\varepsilon=1$ имеем
$y_k(0)=\xi_kD_+(0)+o(1)$, $y_k'(0)/\rho_k=\zeta_k\tau_-(0)D_-(0)+o(1)$, при $\varepsilon=-1$
$y_k(0)=\zeta_kD_-(0)+o(1)$, $y_k'(0)/\rho_k=\xi_k\tau_+(0)D_+(0)+o(1)$, так что
$\abs{y_k(0)}+\abs{y_k'(0)/\rho_k}\ge c$.
\end{theorem}

Для сильно регулярной задачи Дирихле при постоянном $a\in(0,1)$ на серии
$\cos(\rho_k^0\ell_+)=0$ из условия $(\xi_k,\zeta_k)\in\ker\Mm(\rho_k^0)$ следует $\zeta_k=0$, а на
серии $\sin(\rho_k^0\ell_-)=0$ "--- $\xi_k=0$; главные части равны $\cos(\rho_k^0x/\sqrt{1+a})$ и
$\sin(\rho_k^0x/\sqrt{1-a})$ (подробнее "--- в конце \S\,7).

Если $a\in W^2_1$, остаток $R$ в представлении \eqref{eq:fsr} имеет порядок $O(1/\rho)$, а не
$o(1)$, поэтому для сильно регулярной задачи обе асимптотики уточняются.

\begin{theorem}\label{thm:o1k}
Пусть выполнены условия теоремы~\ref{thm:ef} и, кроме того, $a\in W^2_1(-1,1)$. Тогда в
нумерации теоремы~\ref{thm:eig}
\[
 \rho_k=\rho_k^0+O(1/k),\qquad\lambda_k-\lambda_k^0=O(1),
\]
и в \eqref{eq:efas} остатки $o(1)$ можно заменить на $O(1/k)$.
\end{theorem}

\section{Сведение к системе первого порядка}

Система $-Mw''=\lambda w$ имеет четырёхмерное пространство решений, тогда как пара
$w(x)=(y(x),y(-x))$ удовлетворяет двум условиям склейки в нуле. Именно эти условия
оставляют две независимые постоянные, на которые действуют исходные краевые формы
в точках $\pm1$. Поэтому сначала установим соответствие между решениями скалярной
краевой задачи и решениями системы с четырьмя условиями на $[0,1]$.
Пусть
$\Phi(x,\lambda)$ "--- матрица Коши ($\Phi(0,\lambda)=I$) системы
\[
 \mathbf w'=\begin{pmatrix}0&I\\-\lambda M^{-1}(x)&0\end{pmatrix}\mathbf w,\qquad
 \mathbf w=(w,w')\in\C^4,
\]
и
\[
 V_0=\begin{pmatrix}1&-1&0&0\\0&0&1&1\\0&0&0&0\\0&0&0&0\end{pmatrix},\qquad
 V_1=\begin{pmatrix}0&0&0&0\\0&0&0&0\\ \beta_{10}&\alpha_{10}&\beta_{11}&-\alpha_{11}\\
 \beta_{20}&\alpha_{20}&\beta_{21}&-\alpha_{21}\end{pmatrix}.
\]
Строки $V_0$ записывают условия склейки $w_1(0)=w_2(0)$, $w_1'(0)+w_2'(0)=0$, а строки $V_1$ "---
формы $U_j$, выраженные через $\mathbf w(1)$ (формула \eqref{eq:Ufold}). Вместе они означают, что
$\mathbf w(0)$ лежит в ядре матрицы $V_0+V_1\Phi(1,\lambda)$, поэтому положим
$\Delta(\lambda)=\det\bigl(V_0+V_1\Phi(1,\lambda)\bigr)$.

\begin{lemma}\label{lem:fold}
Функция $\Delta$ целая. Число $\lambda$ является собственным значением задачи \eqref{eq:eq},
\eqref{eq:bc} тогда и только тогда, когда $\Delta(\lambda)=0$. При этом отображение
$y\mapsto\mathbf w(0)$, $\mathbf w=(w,w')$, $w(x)=(y(x),y(-x))$, есть изоморфизм пространства
собственных функций на $\ker\bigl(V_0+V_1\Phi(1,\lambda)\bigr)$.
\end{lemma}

\begin{proof}
Пусть $y\in W^2_1(-1,1)$ удовлетворяет \eqref{eq:eq} почти всюду, и $w=(w_1,w_2)$,
$w_1(x)=y(x)$, $w_2(x)=y(-x)$, $x\in[0,1]$. Тогда $w_2'(x)=-y'(-x)$, $w_2''(x)=y''(-x)$, и
уравнение \eqref{eq:eq} в точках $x$ и $-x$ даёт почти всюду на $(0,1)$
\[
 -w_1''-a(x)w_2''=\lambda w_1,\qquad -w_2''-a(-x)w_1''=\lambda w_2,
\]
то есть $-Mw''=\lambda w$. Так как $\det M=1-p\ge1-p_0>0$ и $M^{-1}$ непрерывна, имеем
$w''=-\lambda M^{-1}w$, и $w\in C^2[0,1]$. В нуле $w_1(0)=w_2(0)$ и $w_1'(0)+w_2'(0)=0$; в
единице $y(\pm1)$ и $y'(1)=w_1'(1)$, $y'(-1)=-w_2'(1)$ выражаются через $\mathbf w(1)$, так что
\begin{equation}\label{eq:Ufold}
 U_j(y)=\beta_{j0}w_1(1)+\alpha_{j0}w_2(1)+\beta_{j1}w_1'(1)-\alpha_{j1}w_2'(1).
\end{equation}
Следовательно, $(V_0+V_1\Phi(1,\lambda))\mathbf w(0)=0$. Обратно, если $\mathbf w(0)$ лежит в ядре
этой матрицы и $\mathbf w=\Phi(\cdot,\lambda)\mathbf w(0)$, то функция $y(x)=w_1(x)$ при $x\ge0$,
$y(x)=w_2(-x)$ при $x<0$ принадлежит $C^1[-1,1]$ (значения и односторонние производные в нуле
совпадают), её производная кусочно непрерывно дифференцируема, поэтому $y\in W^2_\infty(-1,1)$;
уравнение \eqref{eq:eq} выполнено на каждой половине отрезка, а условия \eqref{eq:bc} "--- в силу
\eqref{eq:Ufold}. Отображение $y\mapsto\mathbf w(0)$ линейно и инъективно. Коэффициенты системы
непрерывны по $x$ и линейны по $\lambda$, поэтому $\Phi(1,\cdot)$, а с ней и $\Delta$, "--- целые
функции.
\end{proof}

Для постоянной матрицы $M$ собственные векторы $\mathbf u_\pm$ разделяют систему
на два независимых уравнения. При переменном $a$ этот переход точного разделения
не даёт: дифференцирование базиса порождает связь между компонентами. Поэтому
экспоненты $e^{\pm i\rho t_\pm}$ сами по себе ещё не дают асимптотику решений.
Нужно учесть изменение амплитуд и доказать, что оставшаяся поправка стремится к нулю.

Дифференцирование экспонент $e^{\pm i\rho t_\pm}$ даёт множитель $\rho$;
поэтому вместе со значениями решения будем рассматривать производную, делённую на $\rho$.
При $\rho\ne0$ и $\lambda=\rho^2$ положим $Y=(w,w'/\rho)$. Тогда $Y'=\rho KY$, где
$K=\begin{pmatrix}0&I\\-M^{-1}&0\end{pmatrix}$, и векторы
\[
 g_{\pm,s}(x)=\bigl(\mathbf u_\pm(x),\ is\tau_\pm(x)\mathbf u_\pm(x)\bigr),\qquad s=\pm1,
\]
удовлетворяют $Kg_{\pm,s}=is\tau_\pm g_{\pm,s}$. Обозначим через $T(x)$ матрицу со столбцами
$g_{+,+},g_{+,-},g_{-,+},g_{-,-}$; в этом порядке нумеруются столбцы всех матриц ниже.
В этом базисе главная матрица диагональна, а изменение базиса даёт слагаемое
$B=-T^{-1}T'$. Его диагональные элементы определяют амплитуды $D_\pm$, уже выписанные
в п.~\ref{ss:ef}; внедиагональные элементы учитываются в остатке. Соберём амплитуды
и экспоненты в матрицы:
\begin{equation}\label{eq:D}
 \Dd=\operatorname{diag}(D_+,D_+,D_-,D_-),\qquad
 E(x,\rho)=\operatorname{diag}\bigl(e^{i\rho t_+},e^{-i\rho t_+},e^{i\rho t_-},e^{-i\rho t_-}\bigr).
\end{equation}
Для сравнения нулей характеристического определителя нужны оценки не только
на вещественной оси. Обозначим через $\Pi_+^h=\{\rho:\Imm\rho\ge-h\}$ и
$\Pi_-^h=\{\rho:\Imm\rho\le h\}$ полуплоскости, перекрывающиеся в полосе
$\abs{\Imm\rho}\le h$, где $h\ge0$. В каждой из них главные части решений
имеют одни и те же амплитуды и фазы.

\begin{lemma}\label{lem:fsr}
Для каждого $h\ge0$ в каждой из полуплоскостей $\Pi_\pm^h$ существует фундаментальная матрица
системы $Y'=\rho KY$ вида
\begin{equation}\label{eq:fsr}
\begin{gathered}
 Y_a(x,\rho)=T(x)\Dd(x)\bigl(I+R(x,\rho)\bigr)E(x,\rho),\\
 \upsilon(r):=\sup\bigl\{\abs{R(x,\rho)}:x\in[0,1],\ \rho\in\Pi_\pm^h,\ \abs\rho\ge r\bigr\}\to0
\end{gathered}
\end{equation}
при $r\to\infty$.
\end{lemma}

\begin{proof}
Сначала проверим, что переход к собственным векторам допустим при принятой
гладкости коэффициента. Складывая и вычитая столбцы с одинаковым $\tau_\pm$,
получаем
$\det T=4\tau_+\tau_-\bigl(\det(\mathbf u_+,\mathbf u_-)\bigr)^2=16\nu^2\tau_+\tau_-$.
Функции $\sqrt p$, $\tau_\pm$, $\nu$ абсолютно непрерывны: подкоренные выражения
и знаменатели отделены от нуля. Кроме того,
$1-\sqrt{p_0}\le\sigma_-<\sigma_+\le1+\sqrt{p_0}$ и
$m/\norm a_\infty\le\nu^2\le\norm a_\infty/m$.
Поэтому $T$ абсолютно непрерывна, $\det T$ отделён от нуля, а $T^{-1}$ ограничена.
Замена $Y=TZ$ переводит систему $Y'=\rho KY$ в систему
\[
 Z'-B(x)Z=\rho A(x)Z,\qquad B=-T^{-1}T',\qquad
 A=\operatorname{diag}(i\tau_+,-i\tau_+,i\tau_-,-i\tau_-).
\]
Это система (1.1) из~\cite{KS1} с параметром $\rho$ и без члена, убывающего по параметру.
Элементы $A$ непрерывны, элементы $B$ суммируемы, так как $T'\in L_1$, а $T^{-1}$ ограничена.
Разности диагональных элементов $A$ равны $\pm2i\tau_+$, $\pm2i\tau_-$, $\pm i(\tau_--\tau_+)$,
$\pm i(\tau_++\tau_-)$; каждая имеет вид $e^{i\varphi}\varrho(x)$ с постоянным
$\varphi\in\{\pi/2,3\pi/2\}$ и непрерывной функцией $\varrho>0$, поскольку $\tau_->\tau_+>0$.
Значит, выполнено условие (1.10) из~\cite{KS1}. Критические лучи имеют аргументы
$\pi/2-\varphi$, то есть составляют вещественную ось; секторы между ними после
сдвига на $h$ суть $\Pi_\pm^h$. По теореме~1.5 из~\cite{KS1}
(её формулировка приведена в приложении~\ref{app:KS}) в каждой из этих
полуплоскостей существует фундаментальная матрица вида
$\mathcal D(I+R_1)E$ с $R_1\to0$ равномерно по $x\in[0,1]$, где
$\mathcal D=\operatorname{diag}\exp\int_0^xB_{jj}$.

Для вычисления $\mathcal D$ достаточно диагонали $B$. Строка $T^{-1}$,
отвечающая столбцу $g_{\pm,s}$, равна
$\frac14\bigl((1,\pm1/\nu),(is\tau_\pm)^{-1}(1,\pm1/\nu)\bigr)$:
её произведение на этот столбец равно единице, на остальные "--- нулю.
Отсюда соответствующий диагональный элемент $B=-T^{-1}T'$ равен
$-\nu'/(2\nu)-\tau_\pm'/(2\tau_\pm)$ и не зависит от $s$.
Это логарифмическая производная амплитуды $D_\pm$ из п.~\ref{ss:ef},
поэтому $\mathcal D=\Dd\Dd(0)^{-1}$. Постоянная
диагональная матрица $\Dd(0)$ перестановочна с $E$,
поэтому $\mathcal D(I+R_1)E\,\Dd(0)=\Dd(I+R)E$ с $R=\Dd(0)^{-1}R_1\Dd(0)\to0$. Остаётся
положить $Y_a=TZ$.
\end{proof}

\section{Характеристический определитель}

Определение регулярности дано через $F$, то есть через главные члены краевых форм.
Для спектральных выводов нужно установить, что отброшенные члены действительно
малы в характеристическом определителе полной задачи. При этом недостаточно
малости остатка в каждом решении: вне вещественной оси решения растут
экспоненциально. Сравнивать остаток нужно с тем же множителем
$e^{L\abs{\Imm\rho}}$, который отвечает крайним экспонентам $F$.
Искомая связь имеет вид
\[
 \rho^{1-\kappa_1-\kappa_2}\Delta(\rho^2)=c_\ast\bigl(F(\rho)+\eta(\rho)\bigr),\qquad
 \eta(\rho)=o\bigl(e^{L\abs{\Imm\rho}}\bigr)\quad(\abs\rho\to\infty).
\]
В этом представлении (лемма~\ref{lem:det}) учитываются полные формы $U_j$,
в том числе члены нулевого порядка при $\kappa_j=1$.

В переменной $Y=(w,w'/\rho)$ разделим условие склейки с производной на $\rho$,
а форму $U_j$ "--- на $\rho^{\kappa_j}$. Тогда главные члены не содержат степеней
параметра, а члены нулевого порядка в форме порядка~$1$ получают множитель
$\rho^{-1}$. Обозначим матрицы полученных условий через $\widetilde U_0$,
$\widetilde U_1$. Строки $\widetilde U_0$ равны $(1,-1,0,0)$,
$(0,0,1,1)$, $0$, $0$; первые две строки $\widetilde U_1$ нулевые,
а строка $2+j$ при $\kappa_j=1$ равна $U_j^0+\rho^{-1}U_j^1$, где
$U_j^0=(0,0,\beta_{j1},-\alpha_{j1})$,
$U_j^1=(\beta_{j0},\alpha_{j0},0,0)$.
При $\kappa_j=0$ эта строка равна $U_j^0=(\beta_{j0},\alpha_{j0},0,0)$;
в этом случае положим $U_j^1=0$.

Подстановка $Y=Y_a c$ переводит условия в систему на коэффициенты $c\in\C^4$.
Её матрица равна
\[
 \Uf_a(\rho)=\widetilde U_0Y_a(0,\rho)+\widetilde U_1Y_a(1,\rho),
\]
а через $\Uf^0(\rho)$ обозначим ту же матрицу, в которой $R=0$ и $U_j^1=0$. Строки $1$, $2$
матрицы $\Uf^0$ "--- строки $\widetilde U_0T(0)\Dd(0)$, строка $2+j$ "--- $U_j^0T(1)E(1,\rho)$
(напомним, что $\Dd(1)=I$).

Матрица $\Uf^0$ действует на четыре коэффициента при экспонентах, тогда как
матрица $\Mm$ из \eqref{eq:M} "--- на два коэффициента после выполнения условий
склейки. Чтобы сравнить их определители и затем их ядра, выразим эти два
коэффициента в четырёхмерном пространстве. Обозначим через $\mathbf y^0$
пару главных частей из \eqref{eq:Psi}:
\[
 \mathbf y^0(x)=\Bigl(
 \xi D_+\chi_+\mathbf u_++\zeta D_-\chi_-\mathbf u_-,\
 \xi D_+\tau_+\chi_+'\mathbf u_++\zeta D_-\tau_-\chi_-'\mathbf u_-\Bigr)
\]
(все функции берутся в точке $x$). Постоянный ненулевой множитель между
$\det\Mm$ и $F$ обозначим через $J_U$; его значения для трёх видов форм
вычислены в доказательстве следующей леммы.
Для записи постоянной, возникающей из условий склейки, обозначим через
$\tau_\varepsilon(0)$ число $\tau_+(0)$ при $\varepsilon=1$ и $\tau_-(0)$ при $\varepsilon=-1$.

\begin{lemma}\label{lem:kernel}
Координаты решения, составленного из главных частей, в базисе столбцов $T\Dd E$ находятся по
формулам Эйлера: положим
\[
\begin{gathered}
 \mathbf c(\xi,\zeta)=\frac{\xi}{2}(1,1,0,0)+\frac{i\zeta}{2}(0,0,-1,1)\quad(\varepsilon=1),\\
 \mathbf c(\xi,\zeta)=\frac{i\xi}{2}(-1,1,0,0)+\frac{\zeta}{2}(0,0,1,1)\quad(\varepsilon=-1).
\end{gathered}
\]
Тогда при всех $\rho$
\begin{enumerate}
\item[(a)] $T(x)\Dd(x)E(x,\rho)\,\mathbf c(\xi,\zeta)=\mathbf y^0(x)$, $x\in[0,1]$;
\item[(b)] общее ядро первых двух строк матрицы $\Uf^0(\rho)$ есть
 $\{\mathbf c(\xi,\zeta):(\xi,\zeta)\in\C^2\}$, и
 $\Uf^0(\rho)\,\mathbf c(\xi,\zeta)=\bigl(0,0,\Mm(\rho)(\xi,\zeta)^{\top}\bigr)$;
\item[(c)] $\det\Mm(\rho)=J_UF(\rho)$ и
 $\det\Uf^0(\rho)=c_UF(\rho)$, где $c_U=16\varepsilon\tau_\varepsilon(0)D_+(0)D_-(0)J_U\ne0$.
\end{enumerate}
\end{lemma}

\begin{proof}
(a) Пусть $\varepsilon=1$. Сумма столбцов $(+,+)$ и $(+,-)$ матрицы $T\Dd E$ равна
$D_+\bigl(2\mathbf u_+\cos\rho t_+,\,-2\tau_+\mathbf u_+\sin\rho t_+\bigr)$, а разность столбцов
$(-,-)$ и $(-,+)$ равна $D_-\bigl(-2i\mathbf u_-\sin\rho t_-,\,-2i\tau_-\mathbf u_-\cos\rho t_-\bigr)$
(формулы Эйлера). Умножая первую на $\xi/2$, вторую на $i\zeta/2$ и складывая,
получаем вектор с блоками $\xi D_+\mathbf u_+\cos\rho t_++\zeta D_-\mathbf u_-\sin\rho t_-$
и $-\xi D_+\tau_+\mathbf u_+\sin\rho t_++\zeta D_-\tau_-\mathbf u_-\cos\rho t_-$. Так
как $\mathbf u_\pm=(1,\pm\nu)$, это $\mathbf y^0$. Случай $\varepsilon=-1$ проверяется так же.

(b) Поскольку $E(0,\rho)=I$, по (a) первые две строки $\Uf^0\mathbf c$ равны
$\widetilde U_0\mathbf y^0(0)$, то есть разности первых двух его компонент и сумме третьей и
четвёртой компонент. Обе величины равны нулю: при $\varepsilon=1$ имеем $\nu(0)=1$, $\chi_-(0)=0$
и $\chi_+'(0)=0$, при $\varepsilon=-1$ "--- $\nu(0)=-1$, $\chi_+(0)=0$ и $\chi_-'(0)=0$. Строки $\widetilde U_0T(0)\Dd(0)$ линейно независимы, поэтому их общее ядро
двумерно, а отображение $(\xi,\zeta)\mapsto\mathbf c(\xi,\zeta)$ инъективно. Строка $2+j$ вектора
$\Uf^0\mathbf c$ равна $U_j^0\mathbf y^0(1)$, где $D_\pm(1)=1$. При $\kappa_j=0$ это
\[
 \xi(\beta_{j0}+\nu(1)\alpha_{j0})\chi_+(\rho\ell_+)
 +\zeta(\beta_{j0}-\nu(1)\alpha_{j0})\chi_-(\rho\ell_-),
\]
при $\kappa_j=1$ "---
\[
 \xi\tau_+(1)(\beta_{j1}-\nu(1)\alpha_{j1})\chi_+'(\rho\ell_+)
 +\zeta\tau_-(1)(\beta_{j1}+\nu(1)\alpha_{j1})\chi_-'(\rho\ell_-);
\]
это в точности строки \eqref{eq:M}.

(c) Разлагая определитель \eqref{eq:M}, получаем произведение
$\chi_+(\rho\ell_+)\chi_-(\rho\ell_-)$ с коэффициентом
\[
 (\beta_{10}+\nu(1)\alpha_{10})(\beta_{20}-\nu(1)\alpha_{20})
 -(\beta_{10}-\nu(1)\alpha_{10})(\beta_{20}+\nu(1)\alpha_{20})
 =2\nu(1)\bigl(\alpha_{10}\beta_{20}-\beta_{10}\alpha_{20}\bigr)
\]
для вида $(0,0)$ и произведение $\chi_+'(\rho\ell_+)\chi_-'(\rho\ell_-)$ с коэффициентом
\[
 2\tau_+(1)\tau_-(1)\nu(1)\bigl(\beta_{11}\alpha_{21}-\alpha_{11}\beta_{21}\bigr)
\]
для вида $(1,1)$; для вида $(1,0)$ получаем
\[
 \mathcal A\chi_+'(\rho\ell_+)\chi_-(\rho\ell_-)
 -\mathcal B\chi_+(\rho\ell_+)\chi_-'(\rho\ell_-).
\]
Вместе с формулой \eqref{eq:F} это даёт $\det\Mm=J_UF$, где
\[
 J_U=\begin{cases}
 2\nu(1)(\alpha_{10}\beta_{20}-\beta_{10}\alpha_{20}) & \text{для вида }(0,0),\\
 -2\tau_+(1)\tau_-(1)\nu(1)(\beta_{11}\alpha_{21}-\alpha_{11}\beta_{21}) & \text{для вида }(1,1),\\
 -\varepsilon & \text{для вида }(1,0).
 \end{cases}
\]
Для вида $(0,0)$ этот множитель ненулевой в силу независимости форм,
для вида $(1,1)$ "--- в силу нормированности пары, а для вида $(1,0)$
это следует из $\varepsilon\in\{-1,1\}$.

Для второго равенства обозначим через $\mathcal P$ матрицу со столбцами $\mathbf c(1,0)$,
$\mathbf c(0,1)$, $(1,0,0,0)^\top$, $(0,0,1,0)^\top$.
Разлагая $\det\mathcal P$ по второй и четвёртой строкам, получаем
$\det\mathcal P=-i/4$ при обоих значениях $\varepsilon$. По (b) в матрице
$\Uf^0\mathcal P$ первые два столбца имеют нули в строках $1$, $2$ и столбцы матрицы $\Mm$ в
строках $3$, $4$. Строки $1$, $2$ последних двух столбцов "--- это $D_+(0)\widetilde U_0g_{+,+}(0)$
и $D_-(0)\widetilde U_0g_{-,+}(0)$, где $g_{+,+}(0)=(1,\varepsilon,i\tau_+(0),i\varepsilon\tau_+(0))$,
$g_{-,+}(0)=(1,-\varepsilon,i\tau_-(0),-i\varepsilon\tau_-(0))$; без множителей $D_\pm(0)$ они
образуют матрицу
\[
 \begin{pmatrix}0&2\\2i\tau_+(0)&0\end{pmatrix}\quad(\varepsilon=1),\qquad
 \begin{pmatrix}2&0\\0&2i\tau_-(0)\end{pmatrix}\quad(\varepsilon=-1).
\]
Разложение Лапласа по первым двум строкам даёт
\begin{gather*}
 \det\Uf^0\det\mathcal P=-4i\tau_+(0)D_+(0)D_-(0)\det\Mm\quad(\varepsilon=1),\\
 \det\Uf^0\det\mathcal P=4i\tau_-(0)D_+(0)D_-(0)\det\Mm\quad(\varepsilon=-1),
\end{gather*}
то есть $\det\Uf^0=16\varepsilon\tau_\varepsilon(0)D_+(0)D_-(0)\det\Mm$.
\end{proof}

Таким образом, после выполнения условий склейки определитель главных членов
равен $F$ с точностью до ненулевого постоянного множителя. Для полной матрицы
условий связь с исходным характеристическим определителем имеет вид
\begin{equation}\label{eq:detU}
 \Uf_a(\rho)=\Lambda_\rho\bigl(V_0+V_1\Phi(1,\rho^2)\bigr)S_\rho^{-1}Y_a(0,\rho),\qquad
 \det\Uf_a(\rho)=\rho^{1-\kappa_1-\kappa_2}\Delta(\rho^2)\det Y_a(0,\rho),
\end{equation}
где $S_\rho=\operatorname{diag}(1,1,\rho^{-1},\rho^{-1})$ переводит
$\mathbf w$ в $Y$, а
$\Lambda_\rho=\operatorname{diag}(1,\rho^{-1},\rho^{-\kappa_1},\rho^{-\kappa_2})$
нормирует строки условий. Осталось оценить влияние поправки $R$ и членов
$\rho^{-1}U_j^1$ на определитель.

\begin{lemma}\label{lem:det}
Пусть $a$ удовлетворяет условиям \eqref{eq:H}, формы $U_1$, $U_2$ линейно независимы,
а пара $U_1,U_2$ нормирована. Тогда
\begin{equation}\label{eq:main}
 \rho^{1-\kappa_1-\kappa_2}\Delta(\rho^2)=c_\ast\bigl(F(\rho)+\eta(\rho)\bigr),\qquad
 c_\ast=\frac{\varepsilon\tau_\varepsilon(0)J_U}
 {\abs{\nu(1)}\sqrt{\tau_+(1)\tau_-(1)\tau_+(0)\tau_-(0)}},
\end{equation}
где функция $\eta$ голоморфна при $\rho\ne0$ и
$\sup_{\abs\rho\ge r}\abs{\eta(\rho)}e^{-L\abs{\Imm\rho}}\to0$ при $r\to\infty$.
\end{lemma}

\begin{proof}
Матрица Коши в переменной $Y$ равна $S_\rho\Phi(x,\rho^2)S_\rho^{-1}$,
поэтому $Y_a(x,\rho)=S_\rho\Phi(x,\rho^2)S_\rho^{-1}Y_a(0,\rho)$.
Подставляя это равенство в определение $\Uf_a$ и учитывая
$\widetilde U_j=\Lambda_\rho V_jS_\rho^{-1}$, получаем \eqref{eq:detU}.

Возьмём матрицы $Y_a$ леммы~\ref{lem:fsr} при $h=0$; полуплоскости $\Pi^0_\pm$ покрывают $\C$.
Поправка $R$ может зависеть от выбора полуплоскости, но её оценка в обеих
полуплоскостях равномерна. Оценим $\Uf_a-\Uf^0$. Элементы строк $1$, $2$ равны элементам $\widetilde U_0T(0)\Dd(0)R(0,\rho)$
и по модулю не превосходят $C\upsilon(\abs\rho)$. Элемент строки $2+j$ в столбце $(\pm,s)$ равен
$\bigl(U_j^0T(1)R(1,\rho)+\rho^{-1}U_j^1T(1)(I+R(1,\rho))\bigr)_{\pm,s}e^{is\rho\ell_\pm}$,
то есть $(o(1)+O(\rho^{-1}))e^{is\rho\ell_\pm}$ равномерно в $\Pi^0_\pm$; элементы $\Uf^0$ в этой
строке равны ограниченным величинам, умноженным на тот же множитель. Разложим определители
$\Uf_a$ и $\Uf^0$ по первым двум строкам. Каждое слагаемое есть минор строк $1$, $2$
(ограниченный) на минор строк $3$, $4$ по двум оставшимся столбцам; последний содержит множитель
$e^{i\rho\varpi}$, где $\varpi$ "--- сумма показателей $s\ell_\pm$ двух различных столбцов. Столбцы
с одинаковым $\tau_\pm$ имеют противоположные $s$, поэтому $\varpi\in\{0,\pm d,\pm L\}$ и
$\abs{e^{i\rho\varpi}}\le e^{L\abs{\Imm\rho}}$. Отсюда
\begin{equation}\label{eq:eta1}
 \det\Uf_a(\rho)=c_UF(\rho)+\eta_1(\rho),\qquad
 \abs{\eta_1(\rho)}e^{-L\abs{\Imm\rho}}\to0\quad(\abs\rho\to\infty).
\end{equation}
Далее $\det T(0)=16\tau_+(0)\tau_-(0)$, поскольку $\nu(0)^2=1$, и
$\det\Dd(0)=D_+(0)^2D_-(0)^2$, поэтому
$\det Y_a(0,\rho)=16\tau_+(0)\tau_-(0)D_+(0)^2D_-(0)^2(1+o(1))$. По \eqref{eq:detU}
\[
 \rho^{1-\kappa_1-\kappa_2}\Delta(\rho^2)
 =\frac{c_UF(\rho)+\eta_1(\rho)}{16\tau_+(0)\tau_-(0)D_+(0)^2D_-(0)^2(1+o(1))}.
\]
Так как $D_\pm(0)=\bigl(\abs{\nu(1)}\tau_\pm(1)/\tau_\pm(0)\bigr)^{1/2}$, имеем
$c_U/\bigl(16\tau_+(0)\tau_-(0)D_+(0)^2D_-(0)^2\bigr)=c_\ast$. Положим
$\eta=\rho^{1-\kappa_1-\kappa_2}\Delta(\rho^2)/c_\ast-F$; эта
функция голоморфна при $\rho\ne0$, а оценка
$\abs{F(\rho)}\le\bigl(2\abs\theta+\abs{\theta_d}+\abs{\theta_d'}\bigr)e^{L\abs{\Imm\rho}}$
вместе с \eqref{eq:eta1} даёт требуемое убывание
$\eta$. Матрица $Y_a(0,\rho)$ может зависеть от $\rho$ неголоморфно, но в определение $\eta$
она не входит.

\end{proof}

\begin{remark}\label{rem:Sh}
Определение~\ref{def:reg} согласовано с определениями из~\cite[с.~798]{Sh}. Там рассматривается
система $y'-A(x)y=\mu B(x)y$ с вещественной диагональной матрицей $B$; наша система
$Z'-BZ=\rho AZ$ из доказательства леммы~\ref{lem:fsr} приводится к этому виду, если параметром
считать $i\rho$: $\rho A=i\rho\operatorname{diag}(\tau_+,-\tau_+,\tau_-,-\tau_-)$. Если отбросить в
условиях порядка~$1$ члены нулевого порядка, краевые матрицы $\widetilde U_0T(0)$ и
$\widetilde U_1T(1)$ от параметра не зависят. В упрощённой задаче из~\cite{Sh} матрица
$B=-T^{-1}T'$ заменена своей диагональю; фундаментальная матрица упрощённой системы равна
$\Dd E$, а её характеристический определитель "--- $\det\Uf^0(\rho)=c_UF(\rho)$. Коэффициенты при
крайних экспонентах $e^{\pm i\rho L}$ равны $c_U\theta$ и $c_U\theta'$, а в~\cite{Sh} задача
регулярна, если оба коэффициента при крайних экспонентах отличны от нуля; так как
$\abs{\theta'}=\abs\theta$, это равносильно неравенству $\theta\ne0$. Сильная регулярность
в~\cite{Sh} требует, кроме того, асимптотической простоты и отделённости нулей
характеристического определителя упрощённой задачи; при сопоставлении определений это отвечает
простоте и отделённости нулей $F$ с $\Ree\rho\ge0$, кроме конечного числа, то есть сильной
регулярности. Теорема~\ref{thm:eig} доказана для полных условий $U_j=U_j^0+\rho^{-1}U_j^1$.
\end{remark}

\section{Доказательство теоремы~\ref{thm:eig}}

Ненулевые крайние коэффициенты $F$ выполняют две задачи. Вне горизонтальной
полосы одна из крайних экспонент преобладает, поэтому малый остаток из
\eqref{eq:main} не создаёт больших нулей $\Delta(\rho^2)$.
В полосе требуется нижняя оценка $\abs F$ вдали от её нулей: только с такой
оценкой малость $\eta$ позволяет применить теорему Руше.

Однако регулярность не исключает кратных и сближающихся нулей $F$.
Поэтому сравнение на непересекающихся кругах вокруг отдельных нулей здесь
неприменимо. Объединим пересекающиеся круги и сопоставим нули внутри каждой
компоненты с учётом кратностей. Чтобы из такого сопоставления получить
$\rho_k-\rho_k^0\to0$, диаметры компонент должны стремиться к нулю вместе
с радиусами кругов. Это обеспечит равномерная граница числа нулей на отрезках
полосы фиксированной длины, установленная в следующей лемме.

\begin{lemma}\label{lem:F}
Пусть задача регулярна. Тогда:
\begin{enumerate}
\item[(a)] существует $h_0\ge1$, при котором $\abs{F(\rho)}\ge\frac12\abs\theta e^{L\abs{\Imm\rho}}$
 при $\abs{\Imm\rho}\ge h_0$; в частности, $\Zf$ лежит в полосе $\abs{\Imm\rho}<h_0$;
\item[(b)] для каждого $\delta>0$ существует $c_\delta>0$, при котором
 $\abs{F(\rho)}\ge c_\delta e^{L\abs{\Imm\rho}}$, если $\operatorname{dist}(\rho,\Zf)\ge\delta$;
\item[(c)] существует $N_0\ge1$, при котором каждый прямоугольник $[X,X+1]\times[-h_0,h_0]$,
 $X\in\R$, содержит не больше $N_0$ нулей $F$ с учётом кратности;
\item[(d)] множество $\Zf$ бесконечно, и в полуплоскости $\Ree\rho>0$ лежит бесконечно много нулей $F$.
\end{enumerate}
\end{lemma}

\begin{proof}
(a) По определению регулярной задачи $\theta\ne0$ (п.~\ref{ss:model}). Выберем $h_0\ge1$ так, что
$\abs\theta e^{-2Lh_0}+(\abs{\theta_d}+\abs{\theta_d'})e^{-(L-\abs d)h_0}\le\frac12\abs\theta$. При
$\Imm\rho\ge h_0$ слагаемое $\theta'e^{-i\rho L}$ имеет модуль $\abs\theta e^{L\Imm\rho}$, а остальные
слагаемые \eqref{eq:Fexp} в сумме не превосходят $\frac12\abs\theta e^{L\Imm\rho}$. При
$\Imm\rho\le-h_0$ рассуждение то же со слагаемым $\theta e^{i\rho L}$.

(b) Вне полосы $\abs{\Imm\rho}<h_0$ это (a). Пусть, от противного, точки полосы
$\abs{\Imm\rho}\le h_0$ с $\operatorname{dist}(\rho,\Zf)\ge\delta$ таковы, что
$F(\rho_n)\to0$. Запишем $\rho_n=X_n+z_n$, $X_n=\Ree\rho_n$. Переходя к
подпоследовательности, можно считать, что $z_n\to z_\ast$ и $e^{i\omega X_n}\to u_\omega$,
$\abs{u_\omega}=1$, для каждой из частот $\omega\in\{\pm L,\pm d\}$. Тогда функции
$F_n(z)=F(z+X_n)$ сходятся равномерно на компактах к функции
$F_\ast(z)=\theta u_L e^{iLz}+\theta' u_{-L}e^{-iLz}+\theta_d u_d e^{idz}
+\theta_d' u_{-d}e^{-idz}$; она не равна тождественно нулю (коэффициент при $e^{iLz}$ равен
$\theta u_L\ne0$), и $F_\ast(z_\ast)=0$. Но при больших $n$ функция $F_n$ не имеет нулей в круге
$\abs{z-z_\ast}<\delta/2$, и по теореме Гурвица $F_\ast$ не имеет в нём нулей. Противоречие
показывает, что $\abs F$ отделён от нуля на этом множестве, и остаётся заметить, что
$e^{L\abs{\Imm\rho}}\le e^{Lh_0}$ в полосе.

(c) Круг с центром $z_X=X+\frac12+i(h_0+1)$ радиуса $2h_0+2$ содержит прямоугольник. Центр взят
вне полосы $\abs{\Imm\rho}<h_0$, чтобы значение $F$ в нём было оценено снизу: по (a)
$\abs{F(z_X)}\ge\frac12\abs\theta e^{L(h_0+1)}>0$, и по неравенству Йенсена число нулей в круге
$\abs{\rho-z_X}\le2h_0+2$ не
больше $\log\bigl(\max_{\abs{\rho-z_X}=4h_0+4}\abs F/\abs{F(z_X)}\bigr)/\log2$. Из оценки
$\abs F\le Ce^{L\abs{\Imm\rho}}$ эта величина ограничена числом, не зависящим от $X$.

(d) Функция $F$ не равна тождественно нулю по (a) и по оценке
$\abs{F(\rho)}\le Ce^{L\abs\rho}$ имеет порядок роста не выше единицы. Если бы множество $\Zf$
было конечно, то по теореме Адамара о разложении $F(\rho)$ есть произведение многочлена на
$e^{Q(\rho)}$, где $Q$ "--- многочлен степени не выше единицы. Тогда
\[
 \log\abs{F(iX)}=O(\log\abs X)+\Ree Q(iX),
\]
а $\Ree Q(iX)$ есть линейная функция от $X$. По (a) $\log\abs{F(iX)}\ge L\abs X+\log(\abs\theta/2)$
при $\abs X\ge h_0$. Деление на $X$ и переход к пределу дают, что предел отношения
$\Ree Q(iX)/X$ не меньше $L$ при $X\to+\infty$ и не больше $-L$ при $X\to-\infty$, что невозможно
при $L>0$. Значит, $\Zf$ бесконечно. Все нули лежат в полосе (a), и в каждом
ограниченном множестве их конечное число, поэтому $\abs{\Ree\rho}$ на $\Zf$ не ограничен; так как
$\Zf=-\Zf$, бесконечно много нулей лежит в $\Ree\rho>0$.
\end{proof}

\begin{proof}[Доказательство теоремы~\ref{thm:eig}]
Положим $\widetilde\Delta(\rho)=\rho^{1-\kappa_1-\kappa_2}\Delta(\rho^2)$.
По леммам~\ref{lem:det} и~\ref{lem:F}(a)
$\abs{\widetilde\Delta(iX)}\ge\abs{c_\ast}(\frac12\abs\theta-o(1))e^{LX}$
при $X\to+\infty$, поэтому $\Delta\not\equiv0$.

Сначала построим малые контуры, на которых сравнение с $F$ равномерно,
и применим теорему Руше. Это даст локализацию больших нулей и равенство
их числа внутри каждого контура. Затем учтём конечное число оставшихся
нулей и выберем нумерацию; при сильной регулярности каждый далёкий
контур будет содержать один простой нуль.


\step{1} \textit{Кластер состоит не более чем из $2N_0$ кругов и имеет диаметр не больше
$\frac12$.} Положим $\delta_0=1/(8N_0)$ и для $\delta\in(0,\delta_0]$ обозначим через $O_\delta$
объединение кругов $B(z,\delta)$, $z\in\Zf$; компоненты связности $O_\delta$ назовём
$\delta$-кластерами. Если $N$ кругов
с различными центрами образуют связное объединение, то проекции их центров на вещественную ось
лежат на отрезке длины не больше $2\delta N$, который покрывается $2\delta N+1$ отрезками длины $1$;
центры лежат в полосе $\abs{\Imm\rho}<h_0$, и по лемме~\ref{lem:F}(c)
$N\le N_0(2\delta N+1)\le N/4+N_0$, то есть $N\le4N_0/3<2N_0$. Если бы кластер содержал больше
$2N_0$ кругов, в нём нашёлся бы связный набор ровно из $2N_0+1$ кругов. Любые две точки
кластера соединяются цепочкой его кругов, в которой соседние круги пересекаются; цепочку можно
взять без повторов, поэтому в ней не больше $2N_0$ кругов диаметра $2\delta$. Значит, диаметр
кластера не больше $2\delta\cdot2N_0\le\frac12$. Граница кластера состоит из конечного числа
дуг окружностей, и на ней $\operatorname{dist}(\rho,\Zf)\ge\delta$.

\step{2} \textit{Нули $\widetilde\Delta$ с большими $\abs\rho$ лежат в кластерах, и в далёком
кластере их столько же, сколько нулей $F$; отсюда (a).} По лемме~\ref{lem:det} для каждого
$\delta$ найдётся $r_\delta\ge1$, при котором $\abs{\eta(\rho)}<c_\delta e^{L\abs{\Imm\rho}}$ для
$\abs\rho\ge r_\delta$ (постоянная $c_\delta$ из леммы~\ref{lem:F}(b)). Если $\abs\rho\ge r_\delta$
и $\operatorname{dist}(\rho,\Zf)\ge\delta$, то
$\abs{\widetilde\Delta/c_\ast}\ge\abs F-\abs\eta>0$ по лемме~\ref{lem:F}(b). Значит, нули
$\widetilde\Delta$ с $\abs\rho\ge r_{\delta}$ лежат в $O_{\delta}$; при $\delta=\delta_0$ это даёт (a)
с $h=h_0+\delta_0$. Если $\delta$-кластер $\mathcal V$ лежит в
$\{\abs\rho>r_\delta\}$, то на его границе
$\abs{\widetilde\Delta/c_\ast-F}=\abs\eta<c_\delta e^{L\abs{\Imm\rho}}\le\abs F$, и по теореме Руше
функции $\widetilde\Delta$ и $F$ имеют в $\mathcal V$ одинаковое число нулей.

\step{3} \textit{Учёт оставшихся нулей в переменной $\lambda$.}
Сопоставление в далёких кластерах ещё не даёт нумерации всех нулей $\Delta$ и $G$:
нужно проверить, что оставшихся нулей у этих функций поровну. При этом
$\widetilde\Delta$ может иметь особенность в нуле, поэтому теорему Руше
применим к целым функциям $\Delta(\rho^2)$ и $G(\rho^2)$.
Отрезок $[X,X+1]$ пересекают проекции не более $2N_0$ кругов (их центры лежат в
$[X-\delta_0,X+1+\delta_0]$), и суммарная длина проекций не больше $4\delta_0N_0=\frac12$. Поэтому
найдётся $X_0\ge\max(r_{\delta_0},h_0)+1$, при котором прямая $\Ree\rho=X_0$ не пересекает
$O_{\delta_0}$; так как $\Zf=-\Zf$, то же верно для прямой $\Ree\rho=-X_0$. Рассмотрим квадрат
$\mathcal Q=\{\abs{\Ree\rho}<X_0,\ \abs{\Imm\rho}<X_0\}$. На его вертикальных сторонах
$\operatorname{dist}(\rho,\Zf)\ge\delta_0$ по выбору $X_0$, на горизонтальных "--- поскольку
$\abs{\Imm\rho}=X_0\ge h_0+1$, а $\Zf$ лежит в полосе $\abs{\Imm\rho}<h_0$; кроме того,
$\abs\rho\ge X_0>r_{\delta_0}$. По шагу~2 на границе $\mathcal Q$ нет нулей целых функций
$\Delta(\rho^2)=\rho^{\kappa_1+\kappa_2-1}\widetilde\Delta(\rho)$ и
$c_\ast G(\rho^2)=c_\ast\rho^{\kappa_1+\kappa_2-1}F(\rho)$, а модуль их разности
$\abs{c_\ast\rho^{\kappa_1+\kappa_2-1}\eta}$ меньше модуля второй. По теореме Руше эти функции
имеют в $\mathcal Q$ одинаковое число нулей.

Перейдём к переменной $\lambda=\rho^2$. Квадрат $\mathcal Q$ симметричен относительно нуля, поэтому
он совпадает с прообразом множества $\widehat{\mathcal Q}=\{\rho^2:\rho\in\mathcal Q\}$. Если $f$
"--- целая функция, то её нуль $\lambda^0\ne0$ даёт два нуля $\pm\sqrt{\lambda^0}$ функции
$f(\rho^2)$ той же кратности (отображение $\rho\mapsto\rho^2$ локально биголоморфно при
$\rho\ne0$), а нуль $\lambda^0=0$ "--- нуль $\rho=0$ вдвое большей кратности. Значит, число нулей
$f(\rho^2)$ в $\mathcal Q$ вдвое больше числа нулей $f$ в $\widehat{\mathcal Q}$, и функции
$\Delta$ и $G$ имеют в $\widehat{\mathcal Q}$ поровну нулей. Нули $\Delta(\rho^2)$ вне
$\mathcal Q$ лежат по шагу~2 в $\delta_0$-кластерах из $\{\Ree\rho>X_0\}$ или из
$\{\Ree\rho<-X_0\}$; кластеры из $\{\Ree\rho>X_0\}$ назовём правыми. Квадраты точек правых
кластеров в $\widehat{\mathcal Q}$ не попадают, а квадраты точек левых кластеров совпадают с
квадратами точек правых. Итак, нули $\Delta$ вне $\widehat{\mathcal Q}$ "--- это квадраты нулей
$\widetilde\Delta$ из правых кластеров, с теми же кратностями; то же верно для $G$ и $F$.

\step{4} \textit{Нумерация и (b).} При фиксированном радиусе сравнение даёт лишь
фиксированную границу расстояния между сопоставленными нулями. Для остатка $o(1)$
радиусы должны убывать, но для каждого меньшего радиуса оценка Руше начинает
действовать дальше от начала координат. Выберем радиус в каждом кластере с учётом
этого порога. Пусть $\delta_q=\delta_02^{-q}$. Для правого
$\delta_0$-кластера $\mathcal V$ обозначим через $q(\mathcal V)$ наибольшее целое $q\ge0$, для
которого $q\le\inf_{\mathcal V}\Ree\rho$ и $\mathcal V\subset\{\Ree\rho>r_{\delta_q}\}$. Число
$q=0$ этим условиям удовлетворяет, так как
$\inf_{\mathcal V}\Ree\rho\ge X_0>r_{\delta_0}$, а первое условие оставляет конечное число
кандидатов, так что $q(\mathcal V)$ определено. При фиксированном $q$ оба условия выполнены, как
только $\inf_{\mathcal V}\Ree\rho>\max(q,r_{\delta_q})$; поэтому $q(\mathcal V)\to\infty$, когда
$\inf_{\mathcal V}\Ree\rho\to\infty$. Пусть $q=q(\mathcal V)$. Так как $O_{\delta_q}\subset O_{\delta_0}$,
каждый $\delta_q$-кластер, пересекающий $\mathcal V$, содержится в $\mathcal V$ и, значит, в
$\{\Ree\rho>r_{\delta_q}\}$. По шагу~2 все нули $\widetilde\Delta$ из $\mathcal V$ лежат в
$O_{\delta_q}$, и в каждом $\delta_q$-кластере нулей $\widetilde\Delta$ и $F$ поровну. Сопоставим
нули $\widetilde\Delta$ и $F$ внутри каждого $\delta_{q(\mathcal V)}$-кластера, а нули $\Delta$ и $G$
в $\widehat{\mathcal Q}$ "--- произвольно, и занумеруем пары так, чтобы сначала шли пары из
$\widehat{\mathcal Q}$, а затем пары из правых кластеров в порядке неубывания
$\inf_{\mathcal V}\Ree\rho$ (кластеры с равными $\inf_{\mathcal V}\Ree\rho$ "--- в любом порядке;
их конечное число, так как в каждом ограниченном множестве конечно много нулей $F$). По
лемме~\ref{lem:F}(d) и шагу~2 правых кластеров бесконечно много, и
в каждом есть нули $F$ и столько же нулей $\widetilde\Delta$; поэтому обе последовательности
бесконечны. Расстояние между двумя точками $\delta_q$-кластера не больше его диаметра
$4N_0\delta_q$, поэтому $\abs{\rho_k-\rho_k^0}\le4N_0\delta_{q(\mathcal V_k)}\to0$, где
$\mathcal V_k$ "--- правый кластер, содержащий $\rho_k^0$. Это доказывает (b).

\step{5} \textit{Простота при сильной регулярности, то есть (c).} Пусть нули $F$ с
$\Ree\rho\ge0$ отделены числом $\delta_s>0$, а при
$\abs\rho\ge r_s$ просты. Возьмём $\delta<\min(\delta_0,\delta_s/2)$. Правые $\delta$-кластеры с
$\Ree\rho>\max(X_0,r_s,r_\delta)$ "--- одиночные круги с одним простым нулём $F$; по шагу~2 в
каждом из них ровно один простой нуль $\widetilde\Delta$. В этом нуле $\rho\ne0$, множитель
$\rho^{1-\kappa_1-\kappa_2}$ не обращается в нуль, а отображение $\rho\mapsto\rho^2$ локально
биголоморфно, поэтому соответствующий нуль $\Delta$ тоже прост. Если бы ранг
матрицы $V_0+V_1\Phi(1,\lambda)$ в таком нуле был меньше трёх, все её миноры третьего порядка
обращались бы в нуль, и производная определителя, равная сумме произведений производных
элементов на алгебраические дополнения, тоже обращалась бы в нуль. Значит, ядро одномерно, и по
лемме~\ref{lem:fold} собственное подпространство одномерно.
Теорема доказана.
\end{proof}

\section{Примеры}

Задача Дирихле различает два введённых класса уже при постоянном коэффициенте.
Она всегда регулярна, но её функция $F$ распадается на два тригонометрических
множителя. Нули каждого множителя образуют отделённую серию; нули разных
множителей могут совпадать или сближаться. Поэтому регулярность обеспечивает
сопоставление нулей из теоремы~\ref{thm:eig}, а для сильной регулярности
нужно дополнительно проверить взаимное расположение серий.

\begin{proposition}\label{prop:dir}
Задача Дирихле $y(-1)=y(1)=0$ регулярна при любом $a$, удовлетворяющем \eqref{eq:H}. Нули $F$
образуют серии $\pi(n-\frac12)/\ell_+$ и $\pi n/\ell_-$ при $\varepsilon=1$ и серии $\pi n/\ell_+$ и
$\pi(n-\frac12)/\ell_-$ при $\varepsilon=-1$, $n\in\mathbb Z$. Задача сильно регулярна в том и
только том случае, когда $\ell_+/\ell_-=n_+/n_-$ "--- несократимая дробь с нечётным $n_-$ при
$\varepsilon=1$ и с нечётным $n_+$ при $\varepsilon=-1$.
\end{proposition}

\begin{proof}
Формы $U_1=y(-1)$, $U_2=y(1)$ имеют вид $(0,0)$, поэтому задача регулярна. Пусть
$\varepsilon=1$; тогда $F=\cos\rho\ell_+\sin\rho\ell_-$, и нули каждой серии простые, с шагом
$\pi/\ell_\pm$. Нули разных серий совпадают, если $\ell_+/\ell_-=(2n-1)/(2j)$ при некоторых целых
$n$, $j$.

Пусть $\ell_+/\ell_-=n_+/n_-$ несократима и $n_-$ нечётно. Равенство $(2n-1)n_-=2jn_+$
невозможно, так как слева нечётное число, и
\[
 \Bigl\lvert\frac{\pi(n-\frac12)}{\ell_+}-\frac{\pi j}{\ell_-}\Bigr\rvert
 =\frac{\pi\abs{(2n-1)n_--2jn_+}}{2n_-\ell_+}\ge\frac{\pi}{2n_-\ell_+},
\]
поскольку $\ell_-=n_-\ell_+/n_+$. Нули одной серии отстоят друг от друга на $\pi/\ell_+$ или
$\pi/\ell_-=\pi n_+/(n_-\ell_+)$, и оба числа не меньше $\pi/(2n_-\ell_+)$. Поэтому все нули
простые и отделены числом $\pi/(2n_-\ell_+)$. Если $n_-=2n'$ чётно, то $n_+$ нечётно, и при $j=n'(2q-1)$,
$2n-1=n_+(2q-1)$, $q\ge1$, нули совпадают: двукратных нулей бесконечно много. Если отношение
$\ell_+/\ell_-$ иррационально, то
$\abs{\pi(n-\frac12)/\ell_+-\pi j/\ell_-}=(\pi/\ell_+)\abs{j\ell_+/\ell_--(n-\frac12)}$. По
теореме Кронекера дробные части чисел $j\ell_+/\ell_-$, $j\in\mathbb N$, плотны на $[0,1]$, так как
отношение иррационально, и поэтому бесконечно часто попадают в любую
окрестность точки
 $\frac12$; зазор
между сериями сколь угодно мал при сколь угодно больших $j$. При $\varepsilon=-1$ рассуждение то
же с переменой ролей $n_+$ и $n_-$.
\end{proof}

При постоянном $a\in(0,1)$ чётные решения удовлетворяют уравнению
$-(1+a)y''=\lambda y$, а нечётные "--- уравнению $-(1-a)y''=\lambda y$.
Поэтому собственные значения задачи Дирихле равны
\[
 (1+a)\pi^2(n-\tfrac12)^2,\qquad (1-a)\pi^2n^2,\qquad n\ge1,
\]
то есть в точности квадратам ненулевых нулей $F$. Здесь $\ell_\pm=(1\pm a)^{-1/2}$,
и задача сильно регулярна тогда и только тогда, когда число
$\sqrt{(1-a)/(1+a)}$ рационально и знаменатель его несократимой записи нечётен.
Например, при $a=4/5$ это число равно $1/3$.

В смешанном виде $(1,0)$ функция $F$ не обязана распадаться на произведение.
Простоту и отделённость её нулей можно обеспечить иначе: потребовать, чтобы
слагаемое с меньшей частотой было мало на окружностях вокруг нулей
$\cos\rho L$. При $\mathcal A\ne\mathcal B$ положим
$\gamma=\varepsilon(\mathcal A+\mathcal B)/(\mathcal B-\mathcal A)$. По \eqref{eq:Ftrig}
\[
 F(\rho)=\frac{\varepsilon(\mathcal B-\mathcal A)}2\bigl(\cos\rho L+\gamma\cos\rho d\bigr).
\]
Если в каждом из непересекающихся кругов сохраняется ровно один нуль,
то эти нули просты и равномерно отделены. Неравенство \eqref{eq:gamma}
ниже обеспечивает именно такое сравнение, в том числе при комплексном $\gamma$.

\begin{proposition}\label{prop:10}
Пусть формы имеют вид $(1,0)$, $\mathcal A\ne\mathcal B$, и при некотором $r\in(0,\pi/(2L))$
\begin{equation}\label{eq:gamma}
 \abs\gamma\cosh(\abs dr)<\sin(Lr).
\end{equation}
Тогда каждый нуль $F$ лежит в одном из кругов радиуса $r$ с центрами $\pi(n-\frac12)/L$,
$n\in\mathbb Z$, каждый такой круг содержит ровно один нуль, он простой, и задача сильно
регулярна. Если, кроме того, $\gamma\in\R$, то все нули $F$ вещественны. Условие \eqref{eq:gamma}
выполнено при некотором $r$, если $\abs\gamma\cosh(\pi\abs d/(2L))<1$.
\end{proposition}

\begin{proof}
Пусть $f_\gamma(\rho)=\cos\rho L+\gamma\cos\rho d$ и $z_n=\pi(n-\frac12)/L$; тогда
$\cos(z_nL+u)=(-1)^n\sin u$. При $u=u_1+iu_2$, $\abs u\le\pi/2$, имеем
$\abs{\sin u}^2=\sin^2u_1+\sinh^2u_2\ge\sin^2u_1+u_2^2\ge\sin^2\abs u$: функция
$v\mapsto\sin^2\sqrt v-v$ не возрастает на $[0,\abs u^2]$. Поэтому на окружности
$\abs{\rho-z_n}=r$, где $\abs{\Imm\rho}\le r$,
\[
 \abs{\cos\rho L}\ge\sin(Lr)>\abs\gamma\cosh(\abs dr)\ge\abs{\gamma\cos\rho d},
\]
и по теореме Руше в круге ровно один нуль $f_\gamma$, простой. Из \eqref{eq:gamma} следует
$\abs\gamma<1$. На прямых $\Ree\rho=\pi n/L$ имеем
$\abs{\cos\rho L}=\cosh(L\Imm\rho)>\abs\gamma\cosh(d\Imm\rho)\ge\abs{\gamma\cos\rho d}$, так как
$\abs d<L$, а на прямых $\abs{\Imm\rho}=Y_0$ при больших $Y_0$
$\abs{\cos\rho L}\ge\sinh(LY_0)>\abs\gamma\cosh(dY_0)$. По теореме Руше в каждом прямоугольнике
между прямыми $\Ree\rho=\pi(n-1)/L$ и $\Ree\rho=\pi n/L$ функция $f_\gamma$ имеет ровно один нуль,
как $\cos\rho L$, а при больших $\abs{\Imm\rho}$ нулей нет. Значит, все нули $f_\gamma$ лежат по
одному в кругах вокруг $z_n$ и отделены друг от друга числом $\pi/L-2r>0$. Если $\gamma\in\R$, то
$f_\gamma$ вещественна на вещественной оси, и нуль в круге с вещественным центром вместе с
сопряжённым лежит в том же круге; по единственности он вещественен. Наконец, если
$\abs\gamma\cosh(\pi\abs d/(2L))<1$, то разность $\sin(Lr)-\abs\gamma\cosh(\abs dr)$ стремится к
положительному пределу при $r\to\pi/(2L)-0$.
\end{proof}

Условие \eqref{eq:gamma} достаточно, но не необходимо. Например, при вещественном $\gamma$ с
$\abs\gamma<1$ задача сильно регулярна при любых $L$ и $d$. В самом деле, по доказательству
предложения~\ref{prop:10} в каждом прямоугольнике между прямыми $\Ree\rho=\pi(n-1)/L$ и
$\Ree\rho=\pi n/L$ функция $f_\gamma$ имеет ровно один нуль, а для этого использовалось только
$\abs\gamma<1$. Функция $f_\gamma$ вещественна на вещественной оси и в точке $\pi n/L$ равна
$(-1)^n+\gamma\cos(\pi nd/L)$, то есть имеет знак $(-1)^n$ и модуль не меньше $1-\abs\gamma$.
Поэтому этот нуль вещественный, простой и, так как $\abs{f_\gamma'}\le2L$ на вещественной оси,
отстоит от концов интервала не меньше чем на $(1-\abs\gamma)/(2L)$.

Невыполнение регулярности может означать не потерю точности асимптотики,
а отсутствие дискретного множества собственных значений. Для условий следующего
предложения при чётном $a$ оба краевых требования тождественно выполнены на чётных
решениях, и каждое $\lambda\in\C$ становится собственным значением.

\begin{proposition}\label{prop:refl}
Для условий
\begin{equation}\label{eq:refl}
 y'(1)+y'(-1)=0,\qquad y(1)-y(-1)=0
\end{equation}
вида $(1,0)$ имеем $\mathcal A=\tau_+(1)\bigl(1-\nu(1)^2\bigr)$,
$\mathcal B=\tau_-(1)\bigl(1-\nu(1)^2\bigr)$, и
задача \eqref{eq:eq}, \eqref{eq:refl} регулярна тогда и только тогда, когда $a(1)\ne a(-1)$. Если
$a$ чётна, то каждое $\lambda\in\C$ "--- собственное значение задачи \eqref{eq:eq},
\eqref{eq:refl}.
\end{proposition}

\begin{proof}
Здесь $\alpha_{11}=\beta_{11}=1$, $\beta_{20}=1$, $\alpha_{20}=-1$, остальные коэффициенты равны
нулю; вид $(1,0)$ нормирован. Из \eqref{eq:M} получаем $\mathcal A=\tau_+(1)(1-\nu(1)^2)$ и
$\mathcal B=\tau_-(1)(1-\nu(1)^2)$, откуда
$\mathcal A-\mathcal B=(\tau_+(1)-\tau_-(1))(1-\nu(1)^2)$. Так как $\tau_+(1)<\tau_-(1)$, а
$\nu(1)^2=a(-1)/a(1)$, то $\mathcal A\ne\mathcal B$ ровно при $a(1)\ne a(-1)$. Если $a$ чётна, то
решение задачи Коши $-(1+a)y''=\lambda y$, $y(0)=1$, $y'(0)=0$ чётно, удовлетворяет \eqref{eq:eq},
поскольку $y''(-x)=y''(x)$, и условиям \eqref{eq:refl}.
\end{proof}

Для уравнения $-y''=\lambda y$ (случай $a\equiv0$, условиям \eqref{eq:H} не удовлетворяющий) условия
\eqref{eq:refl} нерегулярны по Биркгофу: функция $\cos(\sqrt\lambda\,x)$ удовлетворяет им при
каждом $\lambda$, и множество собственных значений совпадает с $\C$. Для задачи \eqref{eq:eq},
\eqref{eq:refl} с чётным $a$ множество собственных значений по предложению~\ref{prop:refl} то же,
а при $a(1)\ne a(-1)$ функция $\Delta$ не равна
тождественно нулю, и собственные значения образуют дискретное множество с асимптотикой
теоремы~\ref{thm:eig}. При $a(1)=a(-1)$ имеем $\nu(1)=\varepsilon$, и в матрице
\eqref{eq:M} для условий \eqref{eq:refl} один из столбцов тождественно нулевой:
оба условия в главном порядке выполнены на одной и той же главной части решения.
При $a(1)\ne a(-1)$ оба множителя $1-\nu(1)$ и $1+\nu(1)$ ненулевые, а крайний
коэффициент $F$ отличен от нуля по предложению~\ref{prop:refl}. Сильную регулярность задачи
\eqref{eq:eq}, \eqref{eq:refl} мы не исследуем: для неё
$\gamma=\varepsilon\bigl(\tau_+(1)+\tau_-(1)\bigr)/\bigl(\tau_-(1)-\tau_+(1)\bigr)$, $\abs\gamma>1$, и
предложение~\ref{prop:10} неприменимо.

\section{Доказательство теоремы~\ref{thm:ef}}

Теорема~\ref{thm:eig} устанавливает близость $\rho_k$ и $\rho_k^0$, но
коэффициенты $\xi_k,\zeta_k$ в главной части собственной функции определяются
не положением нуля, а ядром матрицы краевых условий. Поэтому для асимптотики
собственных функций требуется сравнить не только определители, но и их ядра.

Матрицы $\Uf_a(\rho_k)$ и $\Uf^0(\rho_k^0)$ отличаются на $o(1)$.
Из одной этой оценки близость порождающих векторов ядер не следует:
если все миноры третьего порядка главной матрицы также стремятся к нулю,
малая поправка может существенно изменить направление ядра.
Сильная регулярность исключает такое вырождение. Простота и равномерная
отделённость нулей $F$ дают оценку $\abs{F'(\rho_k^0)}\ge c>0$;
поскольку $\det\Uf^0=c_UF$, формула дифференцирования определителя
обеспечивает минор третьего порядка, равномерно отделённый от нуля.
Начнём с оценки $F$ вблизи её нулей.

\begin{lemma}\label{lem:simple}
Пусть задача сильно регулярна. Тогда существуют $\delta>0$, $c_1>0$ и $r_F$, такие что для
каждого нуля $\rho^0$ функции $F$ с $\Ree\rho^0>r_F$ других нулей $F$ в круге
$\abs{\rho-\rho^0}\le\delta$ нет и
\begin{equation}\label{eq:simple}
 \abs{F(\rho)}\ge c_1\abs{\rho-\rho^0}e^{L\abs{\Imm\rho}}\quad\text{при}\ \abs{\rho-\rho^0}\le\delta,\qquad
 \abs{F'(\rho^0)}\ge c_1.
\end{equation}
\end{lemma}

\begin{proof}
Пусть нули $F$ с $\Ree\rho\ge0$ отделены числом $\delta_s>0$ и просты при $\abs\rho\ge r_s$.
Возьмём $\delta<\delta_s/2$, $r_F=r_s+2\delta$, и пусть $c_\delta$ из леммы~\ref{lem:F}(b). В круге
$\overline B(\rho^0,\delta)$ нет других нулей $F$, нуль $\rho^0$ простой, и функция
$g(\rho)=F(\rho)/(\rho-\rho^0)$ голоморфна и не обращается в нуль в этом круге. На окружности
$\abs{\rho-\rho^0}=\delta$ расстояние до $\Zf$ не меньше $\delta$, и
$\abs g\ge c_\delta/\delta$. По принципу минимума модуля $\abs g\ge c_\delta/\delta$ во всём круге.
Поскольку в круге $\abs{\Imm\rho}\le h_0+\delta$, получаем \eqref{eq:simple} с
$c_1=(c_\delta/\delta)e^{-L(h_0+\delta)}$; кроме того, $\abs{F'(\rho^0)}=\abs{g(\rho^0)}\ge c_1$.
\end{proof}

\begin{proof}[Доказательство теоремы~\ref{thm:ef}]
Порождающий вектор ядра $\Uf^0(\rho_k^0)$ выберем как ненулевой столбец
присоединённой матрицы. Оценка леммы~\ref{lem:simple} позволит выбрать этот
столбец с нормой, отделённой от нуля независимо от $k$.
Тот же столбец присоединённой матрицы для $\Uf_a(\rho_k)$ будет отличаться
на $o(1)$ и потому останется ненулевым. После этого близость коэффициентов
перенесём на решения равномерно по всему отрезку.

\step{1} \textit{Оценки матриц краевых условий в полосе.} Пусть $h=h_0+1$ и
$\Sigma=\{\abs{\Imm\rho}\le h\}$. Возьмём матрицу $Y_a$
леммы~\ref{lem:fsr} в полуплоскости $\Pi^h_+\supset\Sigma$, так что одна и та же матрица
$\Uf_a$ служит на всей полосе. В $\Sigma$ элементы строк $1$, $2$ матрицы $\Uf^0$ постоянны, а
элементы строк $3$, $4$ равны постоянным, умноженным на $e^{is\rho\ell_\pm}$, причём
$\abs{e^{is\rho\ell_\pm}}\le e^{hL}$. Поэтому при $\rho,\rho'\in\Sigma$
\[
 \abs{\Uf^0_{ij}(\rho)}\le C_0,\qquad \Bigl\lvert\frac{d}{d\rho}\Uf^0_{ij}(\rho)\Bigr\rvert\le C_0,\qquad
 \abs{\Uf^0_{ij}(\rho)-\Uf^0_{ij}(\rho')}\le C_0\abs{\rho-\rho'}.
\]
Разность $\Uf_a-\Uf^0$ оценивается, как в доказательстве леммы~\ref{lem:det}: в строках $1$, $2$
она равна $\widetilde U_0T(0)\Dd(0)R(0,\rho)$, а в строке $2+j$ "--- строке
$U_j^0T(1)R(1,\rho)+\rho^{-1}U_j^1T(1)(I+R(1,\rho))$, умноженной на $E(1,\rho)$, где
$\abs{E(1,\rho)}\le e^{hL}$. Поэтому в $\Sigma$
$\abs{\Uf_a(\rho)-\Uf^0(\rho)}\le C\bigl(\upsilon(\abs\rho)+\abs\rho^{-1}\bigr)$. По
теореме~\ref{thm:eig} при больших $k$ имеем $\rho_k,\rho_k^0\in\Sigma$, $\abs{\rho_k-\rho_k^0}<\delta$
и $\Ree\rho_k^0>r_F$, где $\delta$, $r_F$ из леммы~\ref{lem:simple}. Разность матриц краевых
условий при $\rho=\rho_k$ и при $\rho=\rho_k^0$ обозначим через $\Gamma_k$:
\begin{equation}\label{eq:Ek}
 \Gamma_k=\Uf_a(\rho_k)-\Uf^0(\rho_k^0),\qquad
 \abs{\Gamma_k}\le C\bigl(\upsilon(\abs{\rho_k})+\abs{\rho_k}^{-1}+\abs{\rho_k-\rho_k^0}\bigr)\to0.
\end{equation}

\step{2} \textit{Ядро матрицы $\Uf^0(\rho_k^0)$ одномерно, и его порождающий вектор
отделён от нуля.} По формуле Якоби
$\frac{d}{d\rho}\det\Uf^0=\sum_{i,j}\frac{d\Uf^0_{ij}}{d\rho}\operatorname{cof}_{ij}\Uf^0$, а по
лемме~\ref{lem:kernel}(c) $\det\Uf^0=c_UF$. По лемме~\ref{lem:simple}
$\abs{c_U}c_1\le\abs{c_UF'(\rho_k^0)}\le C_0\sum_{i,j}\abs{\operatorname{cof}_{ij}\Uf^0(\rho_k^0)}$.
Значит, найдётся номер $j_k$, для которого вектор
$v_k^0=\operatorname{adj}\Uf^0(\rho_k^0)\,\mathbf e_{j_k}$ удовлетворяет
$\abs{v_k^0}\ge\abs{c_U}c_1/(16C_0)$; кроме того, $\abs{v_k^0}\le C$. Так как
$\Uf^0\operatorname{adj}\Uf^0=\det\Uf^0\cdot I$ и $F(\rho_k^0)=0$, вектор $v_k^0$ лежит в ядре
$\Uf^0(\rho_k^0)$; ненулевое алгебраическое дополнение показывает, что ранг $\Uf^0(\rho_k^0)$ равен
трём. По лемме~\ref{lem:kernel}(b) $v_k^0=\mathbf c(\xi_k',\zeta_k')$, где
$\Mm(\rho_k^0)(\xi_k',\zeta_k')^\top=0$, и ядро $\Mm(\rho_k^0)$ одномерно: иначе ядро
$\Uf^0(\rho_k^0)$ было бы двумерным. Это (b). Следовательно,
$(\xi_k',\zeta_k')=\mu_k(\xi_k,\zeta_k)$ с числом $\mu_k\ne0$. Отображение $\mathbf c$ линейно,
инъективно и не зависит от $k$, поэтому
$c\abs{(\xi,\zeta)}\le\abs{\mathbf c(\xi,\zeta)}\le C\abs{(\xi,\zeta)}$. Так как
$\abs{(\xi_k,\zeta_k)}=1$, а $\abs{v_k^0}$ ограничен сверху и отделён от нуля, то же верно для
$\abs{\mu_k}$: $c\le\abs{\mu_k}\le C$.

\step{3} \textit{Выбор близкого вектора ядра полной матрицы.}
Положим $v_k=\operatorname{adj}\Uf_a(\rho_k)\,\mathbf e_{j_k}$.
По \eqref{eq:detU} имеем $\det\Uf_a(\rho_k)=0$, поэтому $\Uf_a(\rho_k)v_k=0$.
Элементы присоединённой матрицы "--- многочлены третьей степени от элементов,
ограниченных в совокупности, и по \eqref{eq:Ek}
$\abs{v_k-v_k^0}\le C\abs{\Gamma_k}\to0$.
Нижняя оценка $\abs{v_k^0}\ge c>0$ из шага~2 даёт $v_k\ne0$ при больших $k$.
Таким образом, получен ненулевой вектор ядра полной матрицы, близкий к
выбранному вектору ядра главной матрицы. Простота нуля $\Delta$ и одномерность
собственного подпространства, то есть пункт (a), уже следуют из
теоремы~\ref{thm:eig}(c).

\step{4} \textit{Асимптотика собственной функции.} Решение $Y_k=Y_a(\cdot,\rho_k)v_k$ системы $Y'=\rho_kKY$ удовлетворяет
$\widetilde U_0Y_k(0)+\widetilde U_1Y_k(1)=0$, поэтому $Y_k=(w_k,w_k'/\rho_k)$, где
$w_k(x)=(y_k(x),y_k(-x))$ для собственной функции $y_k$. По \eqref{eq:fsr} и
лемме~\ref{lem:kernel}(a)
\[
 Y_k=T\Dd E(\cdot,\rho_k^0)v_k^0+T\Dd RE(\cdot,\rho_k)v_k+T\Dd E(\cdot,\rho_k)(v_k-v_k^0)
 +T\Dd\bigl(E(\cdot,\rho_k)-E(\cdot,\rho_k^0)\bigr)v_k^0,
\]
и первое слагаемое равно $\mathbf y^0$ при $\rho=\rho_k^0$, $(\xi,\zeta)=(\xi_k',\zeta_k')$. При
$x\in[0,1]$ и $\rho\in\Sigma$ имеем $\abs{E(x,\rho)}\le e^{hL}$ и
$\abs{E(x,\rho_k)-E(x,\rho_k^0)}\le Le^{hL}\abs{\rho_k-\rho_k^0}$; матрица $T\Dd$ ограничена.
Поэтому остальные слагаемые по модулю не превосходят
$C\bigl(\upsilon(\abs{\rho_k})+\abs{\Gamma_k}+\abs{\rho_k-\rho_k^0}\bigr)\to0$ равномерно по $x$.
Деля на $\mu_k$ и вспоминая определение $\mathbf y^0$, а также то, что
$w_k'=(y_k'(x),-y_k'(-x))$, получаем \eqref{eq:efas}. Значения главной части в нуле следуют
из \eqref{eq:Psi}, так как $t_\pm(0)=0$, $\chi_-(0)=\chi_+'(0)=0$ при $\varepsilon=1$ и
$\chi_+(0)=\chi_-'(0)=0$ при $\varepsilon=-1$; числа $D_\pm(0)$ и $\tau_\pm(0)$
положительны и от $k$ не зависят.
Теорема доказана.
\end{proof}

Для задачи Дирихле при $\varepsilon=1$ и $\ell_+/\ell_-=n_+/n_-$ с нечётным $n_-$
(предложение~\ref{prop:dir}) строки матрицы $\Mm$ равны
$(\nu(1)\cos\rho\ell_+,-\nu(1)\sin\rho\ell_-)$ и $(\cos\rho\ell_+,\sin\rho\ell_-)$. На серии
$\cos\rho\ell_+=0$ имеем $\sin\rho\ell_-\ne0$, откуда
$\zeta_k=0$ и $y_k(x)=\xi_kD_+(x)\cos\rho_k^0t_+(x)+o(1)$, $\abs{\xi_k}=1$; на серии
$\sin\rho\ell_-=0$ получаем $\xi_k=0$ и
$y_k(x)=\zeta_kD_-(x)\sin\rho_k^0t_-(x)+o(1)$, $\abs{\zeta_k}=1$, $x\in[0,1]$. При постоянном
$a\in(0,1)$ имеем $\nu\equiv1$, $D_\pm\equiv1$, и главные части равны
$\xi_k\cos(\rho_k^0x/\sqrt{1+a})$ и $\zeta_k\sin(\rho_k^0x/\sqrt{1-a})$ с $\abs{\xi_k}=\abs{\zeta_k}=1$.

\section{Доказательство теоремы~\ref{thm:o1k}}

До сих пор гладкость $a\in W^1_1$ давала лишь стремление остатка к нулю,
без оценки скорости. При $a\in W^2_1$ остаток в решениях имеет порядок
$O(1/\rho)$. Однако скорость приближения собственных значений зависит ещё
и от поведения $F$ вблизи нулей: малая величина $F(\rho_k)$ сама по себе
не даёт такой же оценки расстояния до нуля $F$.

Для сильно регулярной задачи это препятствие снимает
лемма~\ref{lem:simple}: её нижняя оценка линейна по расстоянию до нуля.
Поэтому оценка остатка в решениях переносится сначала на
характеристический определитель, затем на положение его нулей и на
коэффициенты собственных функций. Дополнительная гладкость нужна
только для первого из этих переходов "--- уточнения леммы~\ref{lem:fsr}.

\begin{lemma}\label{lem:fsr2}
Пусть $a$ удовлетворяет условиям \eqref{eq:H} и $a\in W^2_1(-1,1)$. Тогда для каждого $h\ge0$ в
каждой из полуплоскостей $\Pi^h_\pm$ существует фундаментальная матрица системы $Y'=\rho KY$ вида
\[
 Y_a(x,\rho)=T(x)\Dd(x)\bigl(I+R(x,\rho)\bigr)E(x,\rho),\qquad
 \sup_{x\in[0,1]}\abs{R(x,\rho)}\le\frac{C}{\abs\rho}
\]
при всех достаточно больших $\abs\rho$.
\end{lemma}

\begin{proof}
Пространство $W^1_1(0,1)$ "--- алгебра, и если $u\in W^2_1$, а $\psi$ дважды непрерывно
дифференцируема на отрезке значений $u$, то $(\psi\circ u)'=\psi'(u)u'\in W^1_1$, то есть
$\psi\circ u\in W^2_1$. Пространство $W^2_1(0,1)$ вложено в $C^1[0,1]$ и потому тоже является
алгеброй; поэтому $p$, $\tau_\pm$, $1/\tau_\pm$, $\nu$, $1/\nu$ принадлежат $W^2_1(0,1)$ вместе со
своими произведениями и частными, которые встречаются ниже.

Как в доказательстве леммы~\ref{lem:fsr}, замена $Y=TZ$ даёт систему $Z'-BZ=\rho AZ$ с
$B=-T^{-1}T'$; это система (1.1) из~\cite{KS2} с параметром $\rho$ при $m=1$: коэффициенты $B_1$ и
$C$ в ней нулевые. Проверим гипотезы теоремы~1.2 из~\cite{KS2} в той форме, в какой они приведены
в приложении~\ref{app:KS}. Диагональные элементы $A$ равны $\pm i\tau_\pm\in W^1_1$. Каждый элемент
$T^{-1}T'$ "--- сумма
произведений функций $1$, $1/\nu$, $1/\tau_\pm$ из $W^2_1$ на функции $\nu'$, $\tau_\pm'$,
$(\tau_\pm\nu)'$ из $W^1_1$, поэтому все элементы $B$ лежат в $W^1_1$; здесь используется условие
$a\in W^2_1$. Члена $C_1$, убывающего по параметру, нет, поэтому соответствующее условие выполнено.
Разности диагональных элементов $A$ имеют вид $e^{i\varphi}\varrho(x)$ с постоянным
$\varphi\in\{\pi/2,3\pi/2\}$ и $\varrho\in\{2\tau_+,2\tau_-,\tau_++\tau_-,\tau_--\tau_+\}\subset W^1_1$,
$\varrho>0$ при всех $x$; это последняя гипотеза теоремы. По теореме~1.2 из~\cite{KS2}
(приложение~\ref{app:KS}) в
каждом из секторов между критическими лучами, сдвинутых на $h$, то есть в $\Pi_\pm^h$, существует
фундаментальная матрица $\mathcal D\bigl(I+\rho^{-1}R_1+\rho^{-1}\mathcal R\bigr)E$, где, как в
доказательстве леммы~\ref{lem:fsr}, $\mathcal D=\Dd\Dd(0)^{-1}$, а
$\mathcal R(x,\rho)\to0$ при $\abs\rho\to\infty$ равномерно по $x\in[0,1]$. По формуле (1.17)
из~\cite{KS2} матрицу $R_1$ можно взять так, что её внедиагональный элемент $(j,k)$ равен
элементу $(j,k)$ матрицы $\mathcal D^{-1}(B-\operatorname{diag}B)\mathcal D$, делённому на
$A_{kk}-A_{jj}$; эти элементы лежат в $W^1_1\subset C[0,1]$, а
$\abs{A_{kk}-A_{jj}}\ge\min\varrho>0$. Диагональные элементы $R_1$ "--- первообразные суммируемых
функций. Значит, $R_1$ ограничена. Умножив фундаментальную матрицу справа на $\Dd(0)$, как в
доказательстве леммы~\ref{lem:fsr}, получаем вид $\Dd(I+R)E$ с
$R=\rho^{-1}\Dd(0)^{-1}(R_1+\mathcal R)\Dd(0)$. Матрица $R_1$ ограничена, а
$\abs{\mathcal R}\to0$, поэтому $\abs R\le C/\abs\rho$ при больших $\abs\rho$, то есть $R$
удовлетворяет требуемой оценке. Остаётся
положить $Y_a=TZ$.
\end{proof}

\begin{lemma}\label{lem:eta2}
Пусть $a$ удовлетворяет условиям \eqref{eq:H}, $a\in W^2_1(-1,1)$, формы $U_1$, $U_2$ линейно
независимы, а пара $U_1,U_2$ нормирована. Тогда существует $C>0$, при котором функция $\eta$ из \eqref{eq:main}
удовлетворяет оценке
\begin{equation}\label{eq:eta2}
 \abs{\eta(\rho)}\le\frac{C}{\abs\rho}e^{L\abs{\Imm\rho}}
\end{equation}
при всех достаточно больших $\abs\rho$.
\end{lemma}

\begin{proof}
Повторим доказательство леммы~\ref{lem:det} с матрицами леммы~\ref{lem:fsr2}. Элементы
$\Uf_a-\Uf^0$ в строках $1$, $2$ имеют порядок $O(\abs\rho^{-1})$, а в строках $3$, $4$ "---
$O(\abs\rho^{-1})e^{is\rho\ell_\pm}$; поэтому в \eqref{eq:eta1}
$\abs{\eta_1}\le C\abs\rho^{-1}e^{L\abs{\Imm\rho}}$. Далее
$\det Y_a(0,\rho)=\det\bigl(T(0)\Dd(0)\bigr)\bigl(1+O(\abs\rho^{-1})\bigr)$, и по доказательству
леммы~\ref{lem:det} $c_U=c_\ast\det\bigl(T(0)\Dd(0)\bigr)$, так что
\[
 \eta=\frac{c_UF+\eta_1}{c_\ast\det(T(0)\Dd(0))(1+O(\abs\rho^{-1}))}-F
 =\frac{\eta_1/c_U-F\,O(\abs\rho^{-1})}{1+O(\abs\rho^{-1})}
 =O\bigl(\abs\rho^{-1}\bigr)e^{L\abs{\Imm\rho}}.
\]
При больших $\abs\rho$ знаменатель ограничен снизу, поэтому отсюда следует \eqref{eq:eta2}. Функция
$\eta$ определена через $\Delta$ и от выбора $Y_a$ не зависит.
\end{proof}

Оценка \eqref{eq:eta2} регулярности не требует.

\begin{proof}[Доказательство теоремы~\ref{thm:o1k}]
По теореме~\ref{thm:eig} при больших $k$ имеем $\abs{\rho_k-\rho_k^0}<\delta$ и
$\Ree\rho_k^0>r_F$ ($\delta$, $r_F$ из леммы~\ref{lem:simple}). Поскольку $\Delta(\rho_k^2)=0$,
равенство \eqref{eq:main} даёт $F(\rho_k)=-\eta(\rho_k)$, и по \eqref{eq:simple}, \eqref{eq:eta2}
\[
 c_1\abs{\rho_k-\rho_k^0}e^{L\abs{\Imm\rho_k}}\le\abs{F(\rho_k)}=\abs{\eta(\rho_k)}
 \le\frac{C}{\abs{\rho_k}}e^{L\abs{\Imm\rho_k}},
\]
то есть $\abs{\rho_k-\rho_k^0}\le C/(c_1\abs{\rho_k})$. Свяжем $\abs{\rho_k}$ с номером. В нумерации
теоремы~\ref{thm:eig} номерам, меньшим $k$, отвечают нули в
$\widehat{\mathcal Q}$, число которых не зависит от $k$, и нули $F$ из правых кластеров, идущих не позже кластера $\mathcal V_k$,
содержащего $\rho_k^0$. Кластеры упорядочены по $\inf_{\mathcal V}\Ree\rho$, и диаметр каждого не
больше $\frac12$, поэтому эти нули лежат в $\{X_0<\Ree\rho\le\Ree\rho_k^0+1\}$, и по
лемме~\ref{lem:F}(c) их не больше $N_0(\Ree\rho_k^0+2)$. Значит,
$k\le N_0(\Ree\rho_k^0+2)+O(1)$, откуда $\abs{\rho_k}\ge\Ree\rho_k^0-\delta\ge ck$ при
больших $k$, и $\rho_k-\rho_k^0=O(1/k)$. Далее
$\abs{\lambda_k-\lambda_k^0}=\abs{\rho_k-\rho_k^0}\abs{\rho_k+\rho_k^0}\le
C(2\abs{\rho_k}+\delta)/(c_1\abs{\rho_k})=O(1)$.

Наконец, повторим доказательство теоремы~\ref{thm:ef} с матрицей $Y_a$ леммы~\ref{lem:fsr2} в
$\Pi^h_+$. Теперь $\upsilon(r)\le C/r$ и $\abs{\rho_k-\rho_k^0}=O(1/k)$, поэтому в \eqref{eq:Ek}
$\abs{\Gamma_k}=O(1/k)$, и все остатки в шаге~4 имеют порядок $O(1/k)$.
Теорема доказана.
\end{proof}

В статье оценки $O(1/k)$ доказаны для сильно регулярных задач: доказательство опирается на
лемму~\ref{lem:simple}. Для регулярной задачи с кратными или сближающимися нулями $F$ из
\eqref{eq:eta2} следует лишь, что $\abs{F(\rho_k)}=O(1/k)\,e^{L\abs{\Imm\rho_k}}$, и скорость
сходимости $\rho_k-\rho_k^0$ зависит от поведения $F$ вблизи нулей: например, если нули $F$
отделены, но среди них бесконечно много двукратных, то тем же рассуждением получается лишь
$\rho_k-\rho_k^0=O(k^{-1/2})$. Для вида $(1,0)$ при условии \eqref{eq:gamma} и для
задачи Дирихле с несократимым отношением $\ell_+/\ell_-=n_+/n_-$ нужной чётности задача сильно
регулярна, и оценки теоремы~\ref{thm:o1k} выполнены при любом $a\in W^2_1$, удовлетворяющем
\eqref{eq:H}.

\section{Замечания}

1. Условия \eqref{eq:H} требуют, чтобы $a$ была отделена от нуля и $p\le p_0<1$. Если $a$ меняет
знак (то есть нарушено условие (ii)), если $p(x)<0$ в некоторых точках или если $p$ достигает
единицы, числа $\pm i\tau_\pm$ перестают быть различными и чисто мнимыми: при $p<0$
числа $\sigma_\pm$ невещественны, в нулях $p$ совпадают $\tau_+$ и $\tau_-$, а при $p\to1$ функция
$\tau_-$ неограничена. Тогда лемма~\ref{lem:fsr} не применима. Эти случаи требуют иного подхода и
здесь не рассматриваются.

2. В теоремах~\ref{thm:ef} и~\ref{thm:o1k} сильная регулярность используется через
лемму~\ref{lem:simple}: без отделённости нулей $F$ величина $\abs{F'(\rho_k^0)}$ может стремиться к
нулю, и ядро матрицы $\Uf_a(\rho_k)$ не обязано быть близко к ядру $\Uf^0(\rho_k^0)$. Для задачи
Дирихле с иррациональным отношением $\ell_+/\ell_-$ теорема~\ref{thm:eig} даёт
$\rho_k-\rho_k^0\to0$; нам неизвестно, верна ли для неё асимптотика собственных функций вида
\eqref{eq:efas} и какова скорость сходимости $\rho_k-\rho_k^0$. При двукратных нулях $F$ (задача
Дирихле с чётным $n_-$ при $\varepsilon=1$ и с чётным $n_+$ при $\varepsilon=-1$) асимптотика пары
собственных функций здесь не получена.

3. Нумерация в теореме~\ref{thm:eig} учитывает кратности нулей $\Delta$. Совпадение кратности
нуля $\Delta$ с алгебраической кратностью собственного значения оператора $L_a$ здесь не
доказывается; простота в теоремах~\ref{thm:eig}
и~\ref{thm:ef} "--- простота нуля $\Delta$ и одномерность собственного подпространства.

4. В теореме~\ref{thm:ef} собственная функция нормирована условием на коэффициенты главной части:
$\abs{\xi_k}^2+\abs{\zeta_k}^2=1$. Отсюда следует оценка $\max\abs{y_k}=O(1)$. Значение $y_k$ в
отдельной точке может стремиться к нулю: для задачи Дирихле на серии $\sin\rho\ell_-=0$ имеем
$y_k(0)\to0$. Значение и производная в нуле вместе к нулю не стремятся: по последнему утверждению
теоремы~\ref{thm:ef} $\abs{y_k(0)}+\abs{y_k'(0)/\rho_k}\ge c$ при больших $k$. Нормировка в
$L_2(-1,1)$ и вопрос о базисности системы собственных функций здесь не рассматриваются.


5. Уравнение \eqref{eq:eq} не содержит младших членов. В краевых условиях члены нулевого
порядка при формах порядка~$1$ допускаются: они входят в матрицу $\Uf_a$ с множителем
$\rho^{-1}$ и не влияют ни на $F$, ни на главные части собственных функций. Для вида $(1,0)$ при
$\mathcal A=\mathcal B$ задача нерегулярна; этот случай не рассматривается. Множество собственных
значений в нём бывает устроено иначе, чем в теореме~\ref{thm:eig}: по предложению~\ref{prop:refl}
при чётном $a$ задача
с условиями~\eqref{eq:refl} имеет собственным значением каждое $\lambda\in\C$.

\appendix

\section{Теоремы Косарева--Шкаликова}\label{app:KS}

Приводим формулировки двух теорем, на которые опираются леммы~\ref{lem:fsr} и~\ref{lem:fsr2},
в объёме, нужном здесь. Рассматривается система на отрезке $[0,1]$
\[
 y'-B(x)y-C(x,\lambda)y=\lambda A(x)y,\qquad A=\operatorname{diag}(a_1,\dots,a_n),
\]
и обозначается
\[
 \mathcal D(x)=\operatorname{diag}\Bigl(\exp\int_0^xb_{jj}\Bigr),\qquad
 E(x,\lambda)=\operatorname{diag}\Bigl(\exp\Bigl(\lambda\int_0^xa_j\Bigr)\Bigr).
\]
Пусть разности диагональных элементов имеют постоянные аргументы:
$a_j(x)-a_k(x)=e^{i\theta_{jk}}\varrho_{jk}(x)$ с постоянными $\theta_{jk}\in[0,2\pi)$ и
$\varrho_{jk}>0$. Лучи $\arg\lambda=\pi/2-\theta_{jk}$ называются критическими.

\smallskip\noindent\textbf{Теорема~1.5 из~\cite{KS1}} (с.~270--271). Пусть элементы $A$, $B$,
$C(\cdot,\lambda)$ суммируемы, $\norm{c_{ij}(\cdot,\lambda)}_{L_1}\to0$ при $\lambda\to\infty$,
$\varrho_{jk}\in L_1$ и $\varrho_{jk}>0$ почти всюду. Тогда при фиксированном $\kappa\ge0$ в каждом
секторе между соседними критическими лучами, сдвинутом на $\kappa$ вдоль биссектрисы наружу,
существует фундаментальная матрица
\[
 Y(x,\lambda)=\mathcal D(x)\bigl(I+R(x,\lambda)\bigr)E(x,\lambda),
\]
где $R(x,\lambda)\to0$ при $\lambda\to\infty$ в секторе равномерно по $x\in[0,1]$.

\smallskip\noindent\textbf{Теорема~1.2 из~\cite{KS2}} (с.~927--928) при $m=1$. Пусть
$C(x,\lambda)=\lambda^{-1}B_1(x)+\lambda^{-1}C_1(x,\lambda)$, $a_j\in W^1_1$, диагональные элементы $B$
суммируемы, внедиагональные лежат в $W^1_1$, элементы $B_1$ суммируемы,
$\norm{(C_1)_{ij}(\cdot,\lambda)}_{L_1}\to0$, а $\varrho_{jk}\in W^1_1$ и $\varrho_{jk}(x)>0$ при всех
$x$. Тогда в каждом секторе между соседними критическими лучами, сдвинутом на $\kappa\ge0$ вдоль
биссектрисы наружу, существует фундаментальная матрица
\[
 Y(x,\lambda)=\mathcal D(x)\bigl(I+\lambda^{-1}R_1(x)+\lambda^{-1}\mathcal R(x,\lambda)\bigr)E(x,\lambda),
\]
где $\mathcal R(x,\lambda)\to0$ при $\lambda\to\infty$ в секторе равномерно по $x\in[0,1]$, а
матрица $R_1$ от $\lambda$ не зависит. Её можно взять так (формула (1.17) из~\cite{KS2}), что
внедиагональный элемент $(j,k)$ равен элементу $(j,k)$ матрицы
$\mathcal D^{-1}(B-\operatorname{diag}B)\mathcal D$, делённому на $a_k-a_j$, а диагональные
элементы "--- интегралы от $0$ до $x$ суммируемых функций.

\smallskip
В лемме~\ref{lem:fsr} теорема~1.5 применяется с $\lambda=\rho$, $n=4$,
$A=\operatorname{diag}(i\tau_+,-i\tau_+,i\tau_-,-i\tau_-)$, $B=-T^{-1}T'$ и $C=0$; в
лемме~\ref{lem:fsr2} теорема~1.2 "--- с теми же $A$, $B$ и $B_1=C_1=0$.

% ===== Служебный блок: в статью не входит =====
% Таблицы ниже закомментированы и не печатаются. Они оставлены в исходнике для быстрой
% ориентации (в том числе при работе с языковыми моделями): по ним видно, что означает
% символ и где он введён. При изменении обозначений обновляйте таблицы вместе с текстом.
\iffalse
\section{Обозначения}\label{app:not}

\begin{longtable}{>{$}l<{$}>{\raggedright\arraybackslash}p{0.62\textwidth}l}
\multicolumn{1}{l}{Символ} & Смысл & Где \\ \hline\endhead
\varepsilon,\ m,\ p_0 & знак $a$; нижняя граница $\abs a$; верхняя граница $p$ & п.~\ref{ss:H}\\
p(x) & $a(x)a(-x)$ & п.~\ref{ss:H}\\
w(x),\ M(x) & пара $(y(x),y(-x))$ на $[0,1]$ и матрица системы $-Mw''=\lambda w$ & \S\,1, п.~\ref{ss:chan}\\
\sigma_\pm,\ \mathbf u_\pm & собственные значения $M(x)$ и её собственные векторы $(1,\pm\nu)^\top$ & п.~\ref{ss:H}, \ref{ss:chan}\\
\nu & $\varepsilon\sqrt{a(-x)/a(x)}$; $\nu(1)=\varepsilon\sqrt{a(-1)/a(1)}$ & п.~\ref{ss:chan}\\
\tau_\pm & $\sigma_\pm^{-1/2}$; $\tau_\pm(1)$ входит в производную главной части & п.~\ref{ss:chan}\\
t_\pm,\ \ell_\pm & фазовый интеграл $\int_0^x\tau_\pm$ и фазовая длина $t_\pm(1)$ & п.~\ref{ss:chan}\\
L,\ d & $\ell_++\ell_-$ и $\ell_+-\ell_-$, частоты в $F$ & п.~\ref{ss:chan}, \ref{ss:forms}\\
\rho & $\sqrt\lambda$ & п.~\ref{ss:chan}\\
\kappa_j & порядок формы $U_j$; вид пары форм "--- $(\kappa_1,\kappa_2)$ & п.~\ref{ss:forms}\\
\chi_\pm & тригонометрический множитель главной части: $\cos$ или $\sin$ в зависимости от $\varepsilon$ & п.~\ref{ss:forms}\\
\mathcal A,\ \mathcal B & числа, через которые записаны $F$ и регулярность для вида $(1,0)$ & п.~\ref{ss:forms}\\
F & функция $F$: определитель главной части с точностью до множителя & п.~\ref{ss:model}, \ref{ss:forms}\\
\theta,\ \theta',\ \theta_d,\ \theta_d' & коэффициенты $F$ при $e^{i\rho L}$, $e^{-i\rho L}$, $e^{i\rho d}$, $e^{-i\rho d}$ & п.~\ref{ss:model}\\
\Zf,\ G & множество нулей $F$; функция $G(\rho^2)=\rho^{\kappa_1+\kappa_2-1}F(\rho)$ & п.~\ref{ss:model}, \ref{ss:eig}\\
\delta_s,\ r_s & отделённость нулей $F$ и порог их простоты & опр.~\ref{def:reg}\\
\Mm(\rho) & матрица главных членов нормированной пары форм на главных частях \eqref{eq:Psi} & п.~\ref{ss:forms}\\
D_\pm & амплитуда главной части, $D_\pm(1)=1$ & п.~\ref{ss:ef}\\
\Delta & характеристический определитель & п.~\ref{ss:H}, \S\,3\\
\Phi,\ V_0,\ V_1 & матрица Коши системы для $\mathbf w=(w,w')$; условия склейки и краевые формы & \S\,3\\
K,\ Y & матрица системы $Y'=\rho KY$ для $Y=(w,w'/\rho)$ & \S\,3\\
T,\ g_{\pm,s} & матрица собственных векторов $K$ и её столбцы & \S\,3\\
\Dd,\ E & диагональные матрицы амплитуд $D_\pm$ и экспонент $e^{\pm i\rho t_\pm}$ & \S\,3\\
Y_a,\ R,\ \upsilon & фундаментальная матрица, поправка в ней и наибольший модуль поправки при $\abs\rho\ge r$ & лемма~\ref{lem:fsr}\\
\Pi_\pm^h & полуплоскости $\Imm\rho\ge-h$, $\Imm\rho\le h$ & \S\,3\\
\Lambda_\rho,\ S_\rho & нормировка условий и переход от $\mathbf w$ к $Y$ & \S\,4\\
\widetilde U_0,\ \widetilde U_1 & матрицы условий в переменной $Y$ & \S\,4\\
\Uf_a,\ \Uf^0 & матрица условий на фундаментальной матрице и на решениях с главными частями & \S\,4\\
\mathbf c,\ \mathbf y^0 & координаты решения с главными частями в базисе $g_{\pm,s}$ и само решение & \S\,4\\
J_U,\ c_U,\ c_\ast & постоянные множители: $\det\Mm=J_UF$, $\det\Uf^0=c_UF$, основная формула & \S\,4\\
\eta & добавка в основной формуле & лемма~\ref{lem:det}\\
\end{longtable}

Постоянные доказательств и то, от чего они зависят:

\begin{longtable}{>{$}l<{$}>{\raggedright\arraybackslash}p{0.55\textwidth}l}

\multicolumn{1}{l}{Постоянная} & Зависит от & Где \\ \hline\endhead
h_0 & $\theta$, $\theta_d$, $\theta_d'$, $L$, $d$ & лемма~\ref{lem:F}(a)\\
c_\delta & $\delta$ и $F$ & лемма~\ref{lem:F}(b)\\
N_0 & $h_0$, $L$, $\theta$ & лемма~\ref{lem:F}(c)\\
\delta_0 & $N_0$: $\delta_0=1/(8N_0)$ & теорема~\ref{thm:eig}, шаг~1\\
r_\delta & $c_\delta$ и скорость убывания $\eta$ & теорема~\ref{thm:eig}, шаг~2\\
X_0,\ \mathcal Q,\ \widehat{\mathcal Q} & $r_{\delta_0}$, $h_0$ и расположение нулей $F$ & теорема~\ref{thm:eig}, шаг~3\\
\delta,\ c_1,\ r_F & $\delta_s$, $r_s$, $c_\delta$, $h_0$, $L$ & лемма~\ref{lem:simple}\\
C_0 & $h$, $L$ и коэффициенты форм & теорема~\ref{thm:ef}, шаг~1\\
\end{longtable}

\section{Три вида краевых форм}\label{app:types}

\begin{longtable}{>{$}c<{$}>{$}l<{$}>{$}c<{$}>{$}c<{$}l}
\text{Вид} & F(\rho) & \theta & J_U & Регулярна\\ \hline
(0,0) & \tfrac12(\sin\rho L-\varepsilon\sin\rho d) & -i/4 & 2\nu(1)(\alpha_{10}\beta_{20}-\beta_{10}\alpha_{20}) & всегда\\
(1,1) & \tfrac12(\sin\rho L+\varepsilon\sin\rho d) & -i/4 & -2\tau_+(1)\tau_-(1)\nu(1)(\beta_{11}\alpha_{21}-\alpha_{11}\beta_{21}) & всегда\\
(1,0) & \tfrac12\bigl(\varepsilon(\mathcal B-\mathcal A)\cos\rho L+(\mathcal A+\mathcal B)\cos\rho d\bigr) & \tfrac14\varepsilon(\mathcal B-\mathcal A) & -\varepsilon & $\mathcal A\ne\mathcal B$\\
\end{longtable}
\noindent Запись $F$ через $\chi_\pm$ "--- формула \eqref{eq:F}; во всех случаях
$\det\Mm=J_UF$ и $c_U=16\varepsilon\tau_\varepsilon(0)D_+(0)D_-(0)J_U$.
\fi
% ===== Конец служебного блока =====

\begin{thebibliography}{99}
\bibitem{KopSar} A.~Kopzhassarova, A.~Sarsenbi, Basis Properties of Eigenfunctions of
Second-Order Differential Operators with Involution, \emph{Abstract and Applied Analysis},
\textbf{2012} (2012), Article ID 576843, 6~pp.
\bibitem{Sars} A.~Sarsenbi, A.\,M.~Sarsenbi, On Eigenfunctions of the Boundary Value Problems
for Second Order Differential Equations with Involution, \emph{Symmetry}, \textbf{13}:10
(2021), 1972.
\bibitem{KriSar} L.\,V.~Kritskov, A.\,M.~Sarsenbi, Riesz Basis Property of System of Root
Functions of Second-Order Differential Operator with Involution, \emph{Differential Equations},
\textbf{53}:1 (2017), 33--46.
\bibitem{KSadSar} L.\,V.~Kritskov, M.\,A.~Sadybekov, A.\,M.~Sarsenbi, Nonlocal Spectral Problem
for a Second-Order Differential Equation with an Involution, \emph{Bulletin of the Karaganda
University. Mathematics}, \textbf{91}:3 (2018), 53--60.
\bibitem{Bond} N.\,P.~Bondarenko, Inverse spectral problems for functional-differential
operators with involution, \emph{Journal of Differential Equations}, \textbf{318} (2022),
169--186.
\bibitem{BeiSar} A.~Beisebayeva, A.\,M.~Sarsenbi, Asymptotic Behavior of Eigenvalues and
Eigenfunctions in a Non-Self-Adjoint Problem for the Schr\"odinger Operator with Involution,
\emph{AIMS Mathematics}, \textbf{11}:7 (2026), 22052--22068.
\bibitem{Na} М.\,А.~Наймарк, \emph{Линейные дифференциальные операторы}, Наука, М., 1969.
\bibitem{VS} В.\,Е.~Владыкина, А.\,А.~Шкаликов, Регулярные обыкновенные дифференциальные
операторы с инволюцией, \emph{Матем. заметки}, \textbf{106}:5 (2019), 643--659.
\bibitem{Sh} А.\,А.~Шкаликов, Регулярные спектральные задачи гиперболического типа для системы
обыкновенных дифференциальных уравнений первого порядка, \emph{Матем. заметки}, \textbf{110}:5
(2021), 796--800.
\bibitem{GrSh} Ya.\,A.~Granilshchikova, A.\,A.~Shkalikov, Spectral Properties of a
Differential Operator with Involution, \emph{Moscow University Mathematics Bulletin},
\textbf{77}:4 (2022), 204--208.
\bibitem{ShInv} A.\,A.~Shkalikov, Spectral Problem for a Differential Operator with
Involution, \emph{Proceedings of the Steklov Institute of Mathematics},
\textbf{331}:1 (2025), 284--296.
\bibitem{KS1} А.\,П.~Косарев, А.\,А.~Шкаликов, Асимптотические представления решений
$(n\times n)$-систем обыкновенных дифференциальных уравнений с большим параметром,
\emph{Матем. заметки}, \textbf{116}:2 (2024), 266--289.
\bibitem{KS2} А.\,П.~Косарев, А.\,А.~Шкаликов, Асимптотические разложения решений
$(n\times n)$-систем обыкновенных дифференциальных уравнений с большим параметром,
\emph{Матем. заметки}, \textbf{116}:6 (2024), 923--940.
\end{thebibliography}
\end{document}
